% Options for packages loaded elsewhere
% Options for packages loaded elsewhere
\PassOptionsToPackage{unicode}{hyperref}
\PassOptionsToPackage{hyphens}{url}
\PassOptionsToPackage{dvipsnames,svgnames,x11names}{xcolor}
%
\documentclass[
  spanish,
  letterpaper,
  DIV=11,
  numbers=noendperiod]{scrartcl}
\usepackage{xcolor}
\usepackage[margin=2.5cm]{geometry}
\usepackage{amsmath,amssymb}
\setcounter{secnumdepth}{-\maxdimen} % remove section numbering
\usepackage{iftex}
\ifPDFTeX
  \usepackage[T1]{fontenc}
  \usepackage[utf8]{inputenc}
  \usepackage{textcomp} % provide euro and other symbols
\else % if luatex or xetex
  \usepackage{unicode-math} % this also loads fontspec
  \defaultfontfeatures{Scale=MatchLowercase}
  \defaultfontfeatures[\rmfamily]{Ligatures=TeX,Scale=1}
\fi
\usepackage{lmodern}
\ifPDFTeX\else
  % xetex/luatex font selection
\fi
% Use upquote if available, for straight quotes in verbatim environments
\IfFileExists{upquote.sty}{\usepackage{upquote}}{}
\IfFileExists{microtype.sty}{% use microtype if available
  \usepackage[]{microtype}
  \UseMicrotypeSet[protrusion]{basicmath} % disable protrusion for tt fonts
}{}
\makeatletter
\@ifundefined{KOMAClassName}{% if non-KOMA class
  \IfFileExists{parskip.sty}{%
    \usepackage{parskip}
  }{% else
    \setlength{\parindent}{0pt}
    \setlength{\parskip}{6pt plus 2pt minus 1pt}}
}{% if KOMA class
  \KOMAoptions{parskip=half}}
\makeatother
% Make \paragraph and \subparagraph free-standing
\makeatletter
\ifx\paragraph\undefined\else
  \let\oldparagraph\paragraph
  \renewcommand{\paragraph}{
    \@ifstar
      \xxxParagraphStar
      \xxxParagraphNoStar
  }
  \newcommand{\xxxParagraphStar}[1]{\oldparagraph*{#1}\mbox{}}
  \newcommand{\xxxParagraphNoStar}[1]{\oldparagraph{#1}\mbox{}}
\fi
\ifx\subparagraph\undefined\else
  \let\oldsubparagraph\subparagraph
  \renewcommand{\subparagraph}{
    \@ifstar
      \xxxSubParagraphStar
      \xxxSubParagraphNoStar
  }
  \newcommand{\xxxSubParagraphStar}[1]{\oldsubparagraph*{#1}\mbox{}}
  \newcommand{\xxxSubParagraphNoStar}[1]{\oldsubparagraph{#1}\mbox{}}
\fi
\makeatother

\usepackage{color}
\usepackage{fancyvrb}
\newcommand{\VerbBar}{|}
\newcommand{\VERB}{\Verb[commandchars=\\\{\}]}
\DefineVerbatimEnvironment{Highlighting}{Verbatim}{commandchars=\\\{\}}
% Add ',fontsize=\small' for more characters per line
\usepackage{framed}
\definecolor{shadecolor}{RGB}{241,243,245}
\newenvironment{Shaded}{\begin{snugshade}}{\end{snugshade}}
\newcommand{\AlertTok}[1]{\textcolor[rgb]{0.68,0.00,0.00}{#1}}
\newcommand{\AnnotationTok}[1]{\textcolor[rgb]{0.37,0.37,0.37}{#1}}
\newcommand{\AttributeTok}[1]{\textcolor[rgb]{0.40,0.45,0.13}{#1}}
\newcommand{\BaseNTok}[1]{\textcolor[rgb]{0.68,0.00,0.00}{#1}}
\newcommand{\BuiltInTok}[1]{\textcolor[rgb]{0.00,0.23,0.31}{#1}}
\newcommand{\CharTok}[1]{\textcolor[rgb]{0.13,0.47,0.30}{#1}}
\newcommand{\CommentTok}[1]{\textcolor[rgb]{0.37,0.37,0.37}{#1}}
\newcommand{\CommentVarTok}[1]{\textcolor[rgb]{0.37,0.37,0.37}{\textit{#1}}}
\newcommand{\ConstantTok}[1]{\textcolor[rgb]{0.56,0.35,0.01}{#1}}
\newcommand{\ControlFlowTok}[1]{\textcolor[rgb]{0.00,0.23,0.31}{\textbf{#1}}}
\newcommand{\DataTypeTok}[1]{\textcolor[rgb]{0.68,0.00,0.00}{#1}}
\newcommand{\DecValTok}[1]{\textcolor[rgb]{0.68,0.00,0.00}{#1}}
\newcommand{\DocumentationTok}[1]{\textcolor[rgb]{0.37,0.37,0.37}{\textit{#1}}}
\newcommand{\ErrorTok}[1]{\textcolor[rgb]{0.68,0.00,0.00}{#1}}
\newcommand{\ExtensionTok}[1]{\textcolor[rgb]{0.00,0.23,0.31}{#1}}
\newcommand{\FloatTok}[1]{\textcolor[rgb]{0.68,0.00,0.00}{#1}}
\newcommand{\FunctionTok}[1]{\textcolor[rgb]{0.28,0.35,0.67}{#1}}
\newcommand{\ImportTok}[1]{\textcolor[rgb]{0.00,0.46,0.62}{#1}}
\newcommand{\InformationTok}[1]{\textcolor[rgb]{0.37,0.37,0.37}{#1}}
\newcommand{\KeywordTok}[1]{\textcolor[rgb]{0.00,0.23,0.31}{\textbf{#1}}}
\newcommand{\NormalTok}[1]{\textcolor[rgb]{0.00,0.23,0.31}{#1}}
\newcommand{\OperatorTok}[1]{\textcolor[rgb]{0.37,0.37,0.37}{#1}}
\newcommand{\OtherTok}[1]{\textcolor[rgb]{0.00,0.23,0.31}{#1}}
\newcommand{\PreprocessorTok}[1]{\textcolor[rgb]{0.68,0.00,0.00}{#1}}
\newcommand{\RegionMarkerTok}[1]{\textcolor[rgb]{0.00,0.23,0.31}{#1}}
\newcommand{\SpecialCharTok}[1]{\textcolor[rgb]{0.37,0.37,0.37}{#1}}
\newcommand{\SpecialStringTok}[1]{\textcolor[rgb]{0.13,0.47,0.30}{#1}}
\newcommand{\StringTok}[1]{\textcolor[rgb]{0.13,0.47,0.30}{#1}}
\newcommand{\VariableTok}[1]{\textcolor[rgb]{0.07,0.07,0.07}{#1}}
\newcommand{\VerbatimStringTok}[1]{\textcolor[rgb]{0.13,0.47,0.30}{#1}}
\newcommand{\WarningTok}[1]{\textcolor[rgb]{0.37,0.37,0.37}{\textit{#1}}}

\usepackage{longtable,booktabs,array}
\usepackage{calc} % for calculating minipage widths
% Correct order of tables after \paragraph or \subparagraph
\usepackage{etoolbox}
\makeatletter
\patchcmd\longtable{\par}{\if@noskipsec\mbox{}\fi\par}{}{}
\makeatother
% Allow footnotes in longtable head/foot
\IfFileExists{footnotehyper.sty}{\usepackage{footnotehyper}}{\usepackage{footnote}}
\makesavenoteenv{longtable}
\usepackage{graphicx}
\makeatletter
\newsavebox\pandoc@box
\newcommand*\pandocbounded[1]{% scales image to fit in text height/width
  \sbox\pandoc@box{#1}%
  \Gscale@div\@tempa{\textheight}{\dimexpr\ht\pandoc@box+\dp\pandoc@box\relax}%
  \Gscale@div\@tempb{\linewidth}{\wd\pandoc@box}%
  \ifdim\@tempb\p@<\@tempa\p@\let\@tempa\@tempb\fi% select the smaller of both
  \ifdim\@tempa\p@<\p@\scalebox{\@tempa}{\usebox\pandoc@box}%
  \else\usebox{\pandoc@box}%
  \fi%
}
% Set default figure placement to htbp
\def\fps@figure{htbp}
\makeatother



\ifLuaTeX
\usepackage[bidi=basic]{babel}
\else
\usepackage[bidi=default]{babel}
\fi
% get rid of language-specific shorthands (see #6817):
\let\LanguageShortHands\languageshorthands
\def\languageshorthands#1{}


\setlength{\emergencystretch}{3em} % prevent overfull lines

\providecommand{\tightlist}{%
  \setlength{\itemsep}{0pt}\setlength{\parskip}{0pt}}



 


\usepackage{fvextra}
\DefineVerbatimEnvironment{Highlighting}{Verbatim}{breaklines,commandchars=\\\{\}}
\DefineVerbatimEnvironment{OutputCode}{Verbatim}{breaklines,commandchars=\\\{\}}
\KOMAoption{captions}{tableheading}
\makeatletter
\@ifpackageloaded{caption}{}{\usepackage{caption}}
\AtBeginDocument{%
\ifdefined\contentsname
  \renewcommand*\contentsname{Tabla de contenidos}
\else
  \newcommand\contentsname{Tabla de contenidos}
\fi
\ifdefined\listfigurename
  \renewcommand*\listfigurename{Listado de Figuras}
\else
  \newcommand\listfigurename{Listado de Figuras}
\fi
\ifdefined\listtablename
  \renewcommand*\listtablename{Listado de Tablas}
\else
  \newcommand\listtablename{Listado de Tablas}
\fi
\ifdefined\figurename
  \renewcommand*\figurename{Figura}
\else
  \newcommand\figurename{Figura}
\fi
\ifdefined\tablename
  \renewcommand*\tablename{Tabla}
\else
  \newcommand\tablename{Tabla}
\fi
}
\@ifpackageloaded{float}{}{\usepackage{float}}
\floatstyle{ruled}
\@ifundefined{c@chapter}{\newfloat{codelisting}{h}{lop}}{\newfloat{codelisting}{h}{lop}[chapter]}
\floatname{codelisting}{Listado}
\newcommand*\listoflistings{\listof{codelisting}{Listado de Listados}}
\makeatother
\makeatletter
\makeatother
\makeatletter
\@ifpackageloaded{caption}{}{\usepackage{caption}}
\@ifpackageloaded{subcaption}{}{\usepackage{subcaption}}
\makeatother
\usepackage{bookmark}
\IfFileExists{xurl.sty}{\usepackage{xurl}}{} % add URL line breaks if available
\urlstyle{same}
\hypersetup{
  pdftitle={Regresión logística y análisis discriminante},
  pdflang={es},
  colorlinks=true,
  linkcolor={blue},
  filecolor={Maroon},
  citecolor={Blue},
  urlcolor={Blue},
  pdfcreator={LaTeX via pandoc}}


\title{Regresión logística y análisis discriminante}
\usepackage{etoolbox}
\makeatletter
\providecommand{\subtitle}[1]{% add subtitle to \maketitle
  \apptocmd{\@title}{\par {\large #1 \par}}{}{}
}
\makeatother
\subtitle{Issue \#55 --- MA2003B, equipo 7 (sobre los factores comunes)}
\author{}
\date{}
\begin{document}
\maketitle

\renewcommand*\contentsname{Tabla de contenidos}
{
\hypersetup{linkcolor=}
\setcounter{tocdepth}{3}
\tableofcontents
}

\begin{Shaded}
\begin{Highlighting}[]
\FunctionTok{library}\NormalTok{(nanoparquet)}
\FunctionTok{library}\NormalTok{(dplyr)}

\NormalTok{sima\_scores }\OtherTok{\textless{}{-}} \FunctionTok{read\_parquet}\NormalTok{(}\StringTok{"../data/processed/af\_puntajes\_diarios.parquet"}\NormalTok{)}
\NormalTok{sima\_modelado }\OtherTok{\textless{}{-}} \FunctionTok{read\_parquet}\NormalTok{(}\StringTok{"../data/processed/sima\_diario\_modelado.parquet"}\NormalTok{)}

\CommentTok{\# fecha es hora local de pared sin zona horaria (ver CLAUDE.md); nanoparquet la}
\CommentTok{\# lee como UTC y R la despliega en America/Monterrey, corriéndola 6 horas.}
\CommentTok{\# Fijar el atributo, no convertir (as.POSIXct(tz=) sí correría los datos).}
\FunctionTok{attr}\NormalTok{(sima\_scores}\SpecialCharTok{$}\NormalTok{fecha, }\StringTok{"tzone"}\NormalTok{) }\OtherTok{\textless{}{-}} \StringTok{"UTC"}
\FunctionTok{attr}\NormalTok{(sima\_modelado}\SpecialCharTok{$}\NormalTok{fecha, }\StringTok{"tzone"}\NormalTok{) }\OtherTok{\textless{}{-}} \StringTok{"UTC"}

\FunctionTok{summary}\NormalTok{(sima\_scores)}
\CommentTok{\#\textgreater{}       estacion         fecha                         particion    }
\CommentTok{\#\textgreater{}  Length   :17399   Min.   :2021{-}01{-}01 00:00:00   Length   :17399  }
\CommentTok{\#\textgreater{}  N.unique :   15   1st Qu.:2022{-}03{-}22 00:00:00   N.unique :    2  }
\CommentTok{\#\textgreater{}  N.blank  :    0   Median :2023{-}03{-}09 00:00:00   N.blank  :    0  }
\CommentTok{\#\textgreater{}  Min.nchar:    2   Mean   :2023{-}03{-}20 17:11:43   Min.nchar:   10  }
\CommentTok{\#\textgreater{}  Max.nchar:    4   3rd Qu.:2024{-}03{-}05 00:00:00   Max.nchar:   13  }
\CommentTok{\#\textgreater{}                    Max.   :2025{-}06{-}30 00:00:00                    }
\CommentTok{\#\textgreater{}        F1                  F2                  F3           }
\CommentTok{\#\textgreater{}  Min.   :{-}3.536390   Min.   :{-}5.310439   Min.   :{-}3.906081  }
\CommentTok{\#\textgreater{}  1st Qu.:{-}0.846128   1st Qu.:{-}0.618912   1st Qu.:{-}0.802547  }
\CommentTok{\#\textgreater{}  Median :{-}0.042404   Median : 0.024735   Median : 0.052235  }
\CommentTok{\#\textgreater{}  Mean   :{-}0.006957   Mean   : 0.001833   Mean   :{-}0.004547  }
\CommentTok{\#\textgreater{}  3rd Qu.: 0.794218   3rd Qu.: 0.647795   3rd Qu.: 0.850836  }
\CommentTok{\#\textgreater{}  Max.   : 4.202685   Max.   : 9.929991   Max.   : 4.728334}
\FunctionTok{summary}\NormalTok{(sima\_modelado)}
\CommentTok{\#\textgreater{}       estacion         fecha                            zona      }
\CommentTok{\#\textgreater{}  Length   :23931   Min.   :2021{-}01{-}01 00:00:00   Length   :23931  }
\CommentTok{\#\textgreater{}  N.unique :   15   1st Qu.:2022{-}03{-}04 00:00:00   N.unique :    7  }
\CommentTok{\#\textgreater{}  N.blank  :    0   Median :2023{-}04{-}25 00:00:00   N.blank  :    0  }
\CommentTok{\#\textgreater{}  Min.nchar:    2   Mean   :2023{-}04{-}15 06:31:43   Min.nchar:    3  }
\CommentTok{\#\textgreater{}  Max.nchar:    4   3rd Qu.:2024{-}05{-}28 00:00:00   Max.nchar:    8  }
\CommentTok{\#\textgreater{}                    Max.   :2025{-}06{-}30 00:00:00                    }
\CommentTok{\#\textgreater{}                                                                   }
\CommentTok{\#\textgreater{}       msnm            anio           mes             temporada    }
\CommentTok{\#\textgreater{}  Min.   :334.0   Min.   :2021   Min.   : 1.000   Length   :23931  }
\CommentTok{\#\textgreater{}  1st Qu.:432.0   1st Qu.:2022   1st Qu.: 3.000   N.unique :    4  }
\CommentTok{\#\textgreater{}  Median :520.0   Median :2023   Median : 6.000   N.blank  :    0  }
\CommentTok{\#\textgreater{}  Mean   :517.5   Mean   :2023   Mean   : 6.188   Min.nchar:    5  }
\CommentTok{\#\textgreater{}  3rd Qu.:568.0   3rd Qu.:2024   3rd Qu.: 9.000   Max.nchar:    9  }
\CommentTok{\#\textgreater{}  Max.   :702.0   Max.   :2025   Max.   :12.000                    }
\CommentTok{\#\textgreater{}                                                                   }
\CommentTok{\#\textgreater{}       tipo\_dia        festivo       fin\_de\_semana       TOUT\_max    }
\CommentTok{\#\textgreater{}  Length   :23931   Min.   :0.0000   Min.   :0.0000   Min.   :{-}0.55  }
\CommentTok{\#\textgreater{}  N.unique :    4   1st Qu.:0.0000   1st Qu.:0.0000   1st Qu.:24.76  }
\CommentTok{\#\textgreater{}  N.blank  :    0   Median :0.0000   Median :0.0000   Median :30.17  }
\CommentTok{\#\textgreater{}  Min.nchar:    6   Mean   :0.0435   Mean   :0.2862   Mean   :28.84  }
\CommentTok{\#\textgreater{}  Max.nchar:    9   3rd Qu.:0.0000   3rd Qu.:1.0000   3rd Qu.:33.93  }
\CommentTok{\#\textgreater{}                    Max.   :1.0000   Max.   :1.0000   Max.   :47.58  }
\CommentTok{\#\textgreater{}                                                      NAs    :1272   }
\CommentTok{\#\textgreater{}     SR\_acum          NO\_06\_09         NO2\_06\_09        NOX\_06\_09     }
\CommentTok{\#\textgreater{}  Min.   : 0.000   Min.   :  0.500   Min.   :  0.00   Min.   :  0.55  }
\CommentTok{\#\textgreater{}  1st Qu.: 2.214   1st Qu.:  5.725   1st Qu.:  9.65   1st Qu.: 16.68  }
\CommentTok{\#\textgreater{}  Median : 3.590   Median : 11.175   Median : 15.00   Median : 26.85  }
\CommentTok{\#\textgreater{}  Mean   : 3.475   Mean   : 20.855   Mean   : 17.02   Mean   : 37.69  }
\CommentTok{\#\textgreater{}  3rd Qu.: 4.808   3rd Qu.: 24.375   3rd Qu.: 22.20   3rd Qu.: 46.55  }
\CommentTok{\#\textgreater{}  Max.   :23.791   Max.   :379.300   Max.   :119.42   Max.   :443.45  }
\CommentTok{\#\textgreater{}  NAs    :1018     NAs    :1598      NAs    :1624     NAs    :1584    }
\CommentTok{\#\textgreater{}     CO\_06\_09         WSR\_media           WSR\_min          PRS\_media    }
\CommentTok{\#\textgreater{}  Min.   : 0.0000   Min.   :  0.4367   Min.   :  0.100   Min.   :681.3  }
\CommentTok{\#\textgreater{}  1st Qu.: 0.7602   1st Qu.:  6.2698   1st Qu.:  1.700   1st Qu.:709.8  }
\CommentTok{\#\textgreater{}  Median : 1.2900   Median :  8.2625   Median :  2.700   Median :714.4  }
\CommentTok{\#\textgreater{}  Mean   : 1.4283   Mean   :  8.6782   Mean   :  3.154   Mean   :715.3  }
\CommentTok{\#\textgreater{}  3rd Qu.: 1.9225   3rd Qu.: 10.3917   3rd Qu.:  4.100   3rd Qu.:721.6  }
\CommentTok{\#\textgreater{}  Max.   :11.3925   Max.   :147.6317   Max.   :129.300   Max.   :747.9  }
\CommentTok{\#\textgreater{}  NAs    :1554      NAs    :959        NAs    :959       NAs    :751    }
\CommentTok{\#\textgreater{}     RH\_media       RH\_min\_diurna   viento\_u\_media     viento\_v\_media    }
\CommentTok{\#\textgreater{}  Min.   : 0.9874   Min.   : 0.00   Min.   :{-}115.158   Min.   :{-}98.9563  }
\CommentTok{\#\textgreater{}  1st Qu.:47.1250   1st Qu.:25.00   1st Qu.:  {-}7.129   1st Qu.: {-}1.9403  }
\CommentTok{\#\textgreater{}  Median :57.6250   Median :36.00   Median :  {-}4.124   Median :  0.5571  }
\CommentTok{\#\textgreater{}  Mean   :56.6028   Mean   :37.06   Mean   :  {-}4.245   Mean   :  0.2831  }
\CommentTok{\#\textgreater{}  3rd Qu.:67.2083   3rd Qu.:48.00   3rd Qu.:  {-}1.479   3rd Qu.:  2.6366  }
\CommentTok{\#\textgreater{}  Max.   :99.3889   Max.   :95.00   Max.   :  31.490   Max.   : 39.1337  }
\CommentTok{\#\textgreater{}  NAs    :2067      NAs    :1995    NAs    :1387       NAs    :1387      }
\CommentTok{\#\textgreater{}    PM10\_media      PM2.5\_media        SO2\_media           llovio       }
\CommentTok{\#\textgreater{}  Min.   :  4.00   Min.   :  2.182   Min.   : 0.5045   Min.   :0.00000  }
\CommentTok{\#\textgreater{}  1st Qu.: 40.51   1st Qu.: 12.812   1st Qu.: 3.1208   1st Qu.:0.00000  }
\CommentTok{\#\textgreater{}  Median : 54.69   Median : 18.169   Median : 4.1333   Median :0.00000  }
\CommentTok{\#\textgreater{}  Mean   : 60.18   Mean   : 20.524   Mean   : 4.7485   Mean   :0.08185  }
\CommentTok{\#\textgreater{}  3rd Qu.: 73.82   3rd Qu.: 25.602   3rd Qu.: 5.7250   3rd Qu.:0.00000  }
\CommentTok{\#\textgreater{}  Max.   :317.87   Max.   :167.979   Max.   :61.7750   Max.   :1.00000  }
\CommentTok{\#\textgreater{}  NAs    :805      NAs    :5229      NAs    :1720      NAs    :706      }
\CommentTok{\#\textgreater{}       MDA8         ventanas\_validas   O3\_max\_1h      excede\_anio\_1   }
\CommentTok{\#\textgreater{}  Min.   :  2.344   Min.   : 0.00    Min.   :  1.00   Min.   :0.0000  }
\CommentTok{\#\textgreater{}  1st Qu.: 33.375   1st Qu.:24.00    1st Qu.: 39.00   1st Qu.:0.0000  }
\CommentTok{\#\textgreater{}  Median : 43.125   Median :24.00    Median : 53.00   Median :0.0000  }
\CommentTok{\#\textgreater{}  Mean   : 44.610   Mean   :22.87    Mean   : 55.73   Mean   :0.1065  }
\CommentTok{\#\textgreater{}  3rd Qu.: 54.000   3rd Qu.:24.00    3rd Qu.: 69.00   3rd Qu.:0.0000  }
\CommentTok{\#\textgreater{}  Max.   :147.875   Max.   :24.00    Max.   :265.00   Max.   :1.0000  }
\CommentTok{\#\textgreater{}  NAs    :1402                       NAs    :709      NAs    :1402    }
\CommentTok{\#\textgreater{}  excede\_anio\_3    excede\_anio\_5     excede\_1h     }
\CommentTok{\#\textgreater{}  Min.   :0.0000   Min.   :0.000   Min.   :0.0000  }
\CommentTok{\#\textgreater{}  1st Qu.:0.0000   1st Qu.:0.000   1st Qu.:0.0000  }
\CommentTok{\#\textgreater{}  Median :0.0000   Median :0.000   Median :0.0000  }
\CommentTok{\#\textgreater{}  Mean   :0.1604   Mean   :0.307   Mean   :0.0799  }
\CommentTok{\#\textgreater{}  3rd Qu.:0.0000   3rd Qu.:1.000   3rd Qu.:0.0000  }
\CommentTok{\#\textgreater{}  Max.   :1.0000   Max.   :1.000   Max.   :1.0000  }
\CommentTok{\#\textgreater{}  NAs    :1402     NAs    :1402}
\end{Highlighting}
\end{Shaded}

\subsection{Escala de los puntajes}\label{escala-de-los-puntajes}

Los puntajes factoriales \textbf{no tienen varianza unitaria}, y esto
tiene consecuencias que conviene resolver antes de ajustar nada.

\begin{Shaded}
\begin{Highlighting}[]
\NormalTok{tr }\OtherTok{\textless{}{-}}\NormalTok{ sima\_scores}\SpecialCharTok{$}\NormalTok{particion }\SpecialCharTok{==} \StringTok{"entrenamiento"}
\FunctionTok{round}\NormalTok{(}\FunctionTok{sapply}\NormalTok{(sima\_scores[tr, }\FunctionTok{c}\NormalTok{(}\StringTok{"F1"}\NormalTok{, }\StringTok{"F2"}\NormalTok{, }\StringTok{"F3"}\NormalTok{)], sd), }\DecValTok{4}\NormalTok{)}
\CommentTok{\#\textgreater{}     F1     F2     F3 }
\CommentTok{\#\textgreater{} 1.1425 1.1000 1.2048}
\end{Highlighting}
\end{Shaded}

Un coeficiente de regresión se lee ``por unidad del predictor'', así que
con escalas distintas entre ejes las razones de momios \textbf{no son
comparables entre sí}: un eje con mayor dispersión produce un momio por
unidad más chico sin que eso signifique que importa menos.

Esto corrige un error de la versión anterior de este documento, que
trabajaba sobre las puntuaciones del PCA y afirmaba que sus razones de
momios eran por desviación estándar. No lo eran. Una puntuación de
componente principal tiene varianza igual a su eigenvalor, de modo que
las tres desviaciones estándar valían 1.98, 1.72 y 1.23 --- no 1 --- y
los momios de 1.47, 1.58 y 1.22 que están hoy en el reporte son
\textbf{por unidad}, no por desviación estándar, y no son comparables
entre ejes como el texto afirmaba.

La solución es reescalar a varianza unitaria \textbf{con los momentos de
entrenamiento}, nunca con los de 2025:

\begin{Shaded}
\begin{Highlighting}[]
\NormalTok{mu\_f }\OtherTok{\textless{}{-}} \FunctionTok{sapply}\NormalTok{(sima\_scores[tr, }\FunctionTok{c}\NormalTok{(}\StringTok{"F1"}\NormalTok{, }\StringTok{"F2"}\NormalTok{, }\StringTok{"F3"}\NormalTok{)], mean)}
\NormalTok{sd\_f }\OtherTok{\textless{}{-}} \FunctionTok{sapply}\NormalTok{(sima\_scores[tr, }\FunctionTok{c}\NormalTok{(}\StringTok{"F1"}\NormalTok{, }\StringTok{"F2"}\NormalTok{, }\StringTok{"F3"}\NormalTok{)], sd)}

\ControlFlowTok{for}\NormalTok{ (v }\ControlFlowTok{in} \FunctionTok{c}\NormalTok{(}\StringTok{"F1"}\NormalTok{, }\StringTok{"F2"}\NormalTok{, }\StringTok{"F3"}\NormalTok{)) \{}
\NormalTok{    sima\_scores[[v]] }\OtherTok{\textless{}{-}}\NormalTok{ (sima\_scores[[v]] }\SpecialCharTok{{-}}\NormalTok{ mu\_f[[v]]) }\SpecialCharTok{/}\NormalTok{ sd\_f[[v]]}
\NormalTok{\}}

\CommentTok{\# control: entrenamiento en 0 y 1; validación no tiene por qué estarlo}
\FunctionTok{round}\NormalTok{(}\FunctionTok{sapply}\NormalTok{(sima\_scores[tr, }\FunctionTok{c}\NormalTok{(}\StringTok{"F1"}\NormalTok{, }\StringTok{"F2"}\NormalTok{, }\StringTok{"F3"}\NormalTok{)], sd), }\DecValTok{6}\NormalTok{)}
\CommentTok{\#\textgreater{} F1 F2 F3 }
\CommentTok{\#\textgreater{}  1  1  1}
\FunctionTok{round}\NormalTok{(}\FunctionTok{sapply}\NormalTok{(sima\_scores[}\SpecialCharTok{!}\NormalTok{tr, }\FunctionTok{c}\NormalTok{(}\StringTok{"F1"}\NormalTok{, }\StringTok{"F2"}\NormalTok{, }\StringTok{"F3"}\NormalTok{)], sd), }\DecValTok{3}\NormalTok{)}
\CommentTok{\#\textgreater{}    F1    F2    F3 }
\CommentTok{\#\textgreater{} 1.128 0.775 1.042}
\end{Highlighting}
\end{Shaded}

A partir de aquí, cada razón de momios sí es por desviación estándar del
factor y las tres son comparables entre sí.

\subsection{Unión con la matriz de
diseño}\label{uniuxf3n-con-la-matriz-de-diseuxf1o}

Los puntajes factoriales (\texttt{sima\_scores}) solo traen
\texttt{estacion}, \texttt{fecha}, \texttt{particion} y los tres
factores. Zona, calendario y la variable respuesta viven en
\texttt{sima\_modelado}, así que se unen por
\texttt{(estacion,\ fecha)}.

\begin{Shaded}
\begin{Highlighting}[]
\NormalTok{sima\_disc }\OtherTok{\textless{}{-}}\NormalTok{ sima\_scores }\SpecialCharTok{|\textgreater{}}
    \FunctionTok{inner\_join}\NormalTok{(}
\NormalTok{        sima\_modelado }\SpecialCharTok{|\textgreater{}}
            \FunctionTok{select}\NormalTok{(estacion, fecha, zona, tipo\_dia, temporada, excede\_anio\_5),}
        \AttributeTok{by =} \FunctionTok{c}\NormalTok{(}\StringTok{"estacion"}\NormalTok{, }\StringTok{"fecha"}\NormalTok{)}
\NormalTok{    )}

\FunctionTok{stopifnot}\NormalTok{(}\FunctionTok{nrow}\NormalTok{(sima\_disc) }\SpecialCharTok{==} \FunctionTok{nrow}\NormalTok{(sima\_scores))}
\FunctionTok{summary}\NormalTok{(sima\_disc)}
\CommentTok{\#\textgreater{}       estacion         fecha                         particion    }
\CommentTok{\#\textgreater{}  Length   :17399   Min.   :2021{-}01{-}01 00:00:00   Length   :17399  }
\CommentTok{\#\textgreater{}  N.unique :   15   1st Qu.:2022{-}03{-}22 00:00:00   N.unique :    2  }
\CommentTok{\#\textgreater{}  N.blank  :    0   Median :2023{-}03{-}09 00:00:00   N.blank  :    0  }
\CommentTok{\#\textgreater{}  Min.nchar:    2   Mean   :2023{-}03{-}20 17:11:43   Min.nchar:   10  }
\CommentTok{\#\textgreater{}  Max.nchar:    4   3rd Qu.:2024{-}03{-}05 00:00:00   Max.nchar:   13  }
\CommentTok{\#\textgreater{}                    Max.   :2025{-}06{-}30 00:00:00                    }
\CommentTok{\#\textgreater{}                                                                   }
\CommentTok{\#\textgreater{}        F1                  F2                  F3                   zona      }
\CommentTok{\#\textgreater{}  Min.   :{-}3.095212   Min.   :{-}4.827839   Min.   :{-}3.242120   Length   :17399  }
\CommentTok{\#\textgreater{}  1st Qu.:{-}0.740571   1st Qu.:{-}0.562667   1st Qu.:{-}0.666129   N.unique :    7  }
\CommentTok{\#\textgreater{}  Median :{-}0.037114   Median : 0.022487   Median : 0.043356   N.blank  :    0  }
\CommentTok{\#\textgreater{}  Mean   :{-}0.006089   Mean   : 0.001667   Mean   :{-}0.003774   Min.nchar:    3  }
\CommentTok{\#\textgreater{}  3rd Qu.: 0.695136   3rd Qu.: 0.588925   3rd Qu.: 0.706210   Max.nchar:    8  }
\CommentTok{\#\textgreater{}  Max.   : 3.678384   Max.   : 9.027577   Max.   : 3.924606                    }
\CommentTok{\#\textgreater{}                                                                               }
\CommentTok{\#\textgreater{}       tipo\_dia         temporada     excede\_anio\_5   }
\CommentTok{\#\textgreater{}  Length   :17399   Length   :17399   Min.   :0.0000  }
\CommentTok{\#\textgreater{}  N.unique :    4   N.unique :    4   1st Qu.:0.0000  }
\CommentTok{\#\textgreater{}  N.blank  :    0   N.blank  :    0   Median :0.0000  }
\CommentTok{\#\textgreater{}  Min.nchar:    6   Min.nchar:    5   Mean   :0.3035  }
\CommentTok{\#\textgreater{}  Max.nchar:    9   Max.nchar:    9   3rd Qu.:1.0000  }
\CommentTok{\#\textgreater{}                                      Max.   :1.0000  }
\CommentTok{\#\textgreater{}                                      NAs    :370}
\end{Highlighting}
\end{Shaded}

\subsection{Categóricas: referencia explícita y dummies
k−1}\label{categuxf3ricas-referencia-expluxedcita-y-dummies-k1}

Referencia por variable (issue \#55):

\begin{itemize}
\tightlist
\item
  \texttt{zona}: \textbf{Noreste}, la de menor tasa de excedencia (19.0
  \%, \texttt{docs/05\_estacionalidad\_ozono.pdf}) --- baseline de menor
  riesgo.
\item
  \texttt{tipo\_dia}: \textbf{laborable}, el nivel más frecuente y el
  orden natural de \texttt{NIVELES\_TIPO\_DIA} en
  \texttt{src/agregacion\_diaria.py}.
\item
  \texttt{temporada}: \textbf{invierno}, temporada de menor actividad
  fotoquímica.
\end{itemize}

\texttt{festivo} no se agrega aparte: por construcción
(\texttt{src/agregacion\_diaria.py}, \texttt{marcar\_calendario})
\texttt{tipo\_dia\ ==\ "festivo"} es idéntico a \texttt{festivo\ ==\ 1}
--- son la misma indicadora. Incluir ambas duplicaría la columna y
generaría colinealidad perfecta; la dummy \texttt{tipo\_dia\_festivo} ya
la cubre.

\begin{Shaded}
\begin{Highlighting}[]
\NormalTok{sima\_disc }\OtherTok{\textless{}{-}}\NormalTok{ sima\_disc }\SpecialCharTok{|\textgreater{}}
    \FunctionTok{mutate}\NormalTok{(}
        \AttributeTok{zona =} \FunctionTok{relevel}\NormalTok{(}\FunctionTok{factor}\NormalTok{(zona), }\AttributeTok{ref =} \StringTok{"Noreste"}\NormalTok{),}
        \AttributeTok{tipo\_dia =} \FunctionTok{factor}\NormalTok{(tipo\_dia,}
            \AttributeTok{levels =} \FunctionTok{c}\NormalTok{(}\StringTok{"laborable"}\NormalTok{, }\StringTok{"sabado"}\NormalTok{, }\StringTok{"domingo"}\NormalTok{, }\StringTok{"festivo"}\NormalTok{)}
\NormalTok{        ),}
        \AttributeTok{temporada =} \FunctionTok{factor}\NormalTok{(temporada,}
            \AttributeTok{levels =} \FunctionTok{c}\NormalTok{(}\StringTok{"invierno"}\NormalTok{, }\StringTok{"primavera"}\NormalTok{, }\StringTok{"verano"}\NormalTok{, }\StringTok{"otoño"}\NormalTok{)}
\NormalTok{        )}
\NormalTok{    )}

\NormalTok{dummies }\OtherTok{\textless{}{-}} \FunctionTok{model.matrix}\NormalTok{(}\SpecialCharTok{\textasciitilde{}}\NormalTok{ zona }\SpecialCharTok{+}\NormalTok{ tipo\_dia }\SpecialCharTok{+}\NormalTok{ temporada, }\AttributeTok{data =}\NormalTok{ sima\_disc)[, }\SpecialCharTok{{-}}\DecValTok{1}\NormalTok{]}
\FunctionTok{colnames}\NormalTok{(dummies) }\OtherTok{\textless{}{-}} \FunctionTok{gsub}\NormalTok{(}\StringTok{"\^{}zona|\^{}tipo\_dia|\^{}temporada"}\NormalTok{, }\StringTok{""}\NormalTok{, }\FunctionTok{colnames}\NormalTok{(dummies))}
\FunctionTok{colnames}\NormalTok{(dummies) }\OtherTok{\textless{}{-}} \FunctionTok{c}\NormalTok{(}
    \FunctionTok{paste0}\NormalTok{(}\StringTok{"zona\_"}\NormalTok{, }\FunctionTok{colnames}\NormalTok{(dummies)[}\DecValTok{1}\SpecialCharTok{:}\DecValTok{6}\NormalTok{]),}
    \FunctionTok{paste0}\NormalTok{(}\StringTok{"tipo\_dia\_"}\NormalTok{, }\FunctionTok{colnames}\NormalTok{(dummies)[}\DecValTok{7}\SpecialCharTok{:}\DecValTok{9}\NormalTok{]),}
    \FunctionTok{paste0}\NormalTok{(}\StringTok{"temporada\_"}\NormalTok{, }\FunctionTok{colnames}\NormalTok{(dummies)[}\DecValTok{10}\SpecialCharTok{:}\DecValTok{12}\NormalTok{])}
\NormalTok{)}

\NormalTok{sima\_disc }\OtherTok{\textless{}{-}} \FunctionTok{cbind}\NormalTok{(sima\_disc, dummies)}
\FunctionTok{stopifnot}\NormalTok{(}\FunctionTok{ncol}\NormalTok{(dummies) }\SpecialCharTok{==} \DecValTok{12}\NormalTok{) }\CommentTok{\# 6 zona + 3 tipo\_dia + 3 temporada}
\FunctionTok{summary}\NormalTok{(sima\_disc)}
\CommentTok{\#\textgreater{}       estacion         fecha                         particion    }
\CommentTok{\#\textgreater{}  Length   :17399   Min.   :2021{-}01{-}01 00:00:00   Length   :17399  }
\CommentTok{\#\textgreater{}  N.unique :   15   1st Qu.:2022{-}03{-}22 00:00:00   N.unique :    2  }
\CommentTok{\#\textgreater{}  N.blank  :    0   Median :2023{-}03{-}09 00:00:00   N.blank  :    0  }
\CommentTok{\#\textgreater{}  Min.nchar:    2   Mean   :2023{-}03{-}20 17:11:43   Min.nchar:   10  }
\CommentTok{\#\textgreater{}  Max.nchar:    4   3rd Qu.:2024{-}03{-}05 00:00:00   Max.nchar:   13  }
\CommentTok{\#\textgreater{}                    Max.   :2025{-}06{-}30 00:00:00                    }
\CommentTok{\#\textgreater{}                                                                   }
\CommentTok{\#\textgreater{}        F1                  F2                  F3                  zona     }
\CommentTok{\#\textgreater{}  Min.   :{-}3.095212   Min.   :{-}4.827839   Min.   :{-}3.242120   Noreste :3160  }
\CommentTok{\#\textgreater{}  1st Qu.:{-}0.740571   1st Qu.:{-}0.562667   1st Qu.:{-}0.666129   Centro  :1494  }
\CommentTok{\#\textgreater{}  Median :{-}0.037114   Median : 0.022487   Median : 0.043356   Noroeste:2407  }
\CommentTok{\#\textgreater{}  Mean   :{-}0.006089   Mean   : 0.001667   Mean   :{-}0.003774   Norte   :2080  }
\CommentTok{\#\textgreater{}  3rd Qu.: 0.695136   3rd Qu.: 0.588925   3rd Qu.: 0.706210   Sur     :1347  }
\CommentTok{\#\textgreater{}  Max.   : 3.678384   Max.   : 9.027577   Max.   : 3.924606   Sureste :4114  }
\CommentTok{\#\textgreater{}                                                              Suroeste:2797  }
\CommentTok{\#\textgreater{}       tipo\_dia         temporada    excede\_anio\_5     zona\_Centro     }
\CommentTok{\#\textgreater{}  laborable:11822   invierno :4291   Min.   :0.0000   Min.   :0.00000  }
\CommentTok{\#\textgreater{}  sabado   : 2421   primavera:4822   1st Qu.:0.0000   1st Qu.:0.00000  }
\CommentTok{\#\textgreater{}  domingo  : 2393   verano   :4196   Median :0.0000   Median :0.00000  }
\CommentTok{\#\textgreater{}  festivo  :  763   otoño    :4090   Mean   :0.3035   Mean   :0.08587  }
\CommentTok{\#\textgreater{}                                     3rd Qu.:1.0000   3rd Qu.:0.00000  }
\CommentTok{\#\textgreater{}                                     Max.   :1.0000   Max.   :1.00000  }
\CommentTok{\#\textgreater{}                                     NAs    :370                       }
\CommentTok{\#\textgreater{}  zona\_Noroeste      zona\_Norte        zona\_Sur        zona\_Sureste   }
\CommentTok{\#\textgreater{}  Min.   :0.0000   Min.   :0.0000   Min.   :0.00000   Min.   :0.0000  }
\CommentTok{\#\textgreater{}  1st Qu.:0.0000   1st Qu.:0.0000   1st Qu.:0.00000   1st Qu.:0.0000  }
\CommentTok{\#\textgreater{}  Median :0.0000   Median :0.0000   Median :0.00000   Median :0.0000  }
\CommentTok{\#\textgreater{}  Mean   :0.1383   Mean   :0.1195   Mean   :0.07742   Mean   :0.2365  }
\CommentTok{\#\textgreater{}  3rd Qu.:0.0000   3rd Qu.:0.0000   3rd Qu.:0.00000   3rd Qu.:0.0000  }
\CommentTok{\#\textgreater{}  Max.   :1.0000   Max.   :1.0000   Max.   :1.00000   Max.   :1.0000  }
\CommentTok{\#\textgreater{}                                                                      }
\CommentTok{\#\textgreater{}  zona\_Suroeste    tipo\_dia\_sabado  tipo\_dia\_domingo tipo\_dia\_festivo }
\CommentTok{\#\textgreater{}  Min.   :0.0000   Min.   :0.0000   Min.   :0.0000   Min.   :0.00000  }
\CommentTok{\#\textgreater{}  1st Qu.:0.0000   1st Qu.:0.0000   1st Qu.:0.0000   1st Qu.:0.00000  }
\CommentTok{\#\textgreater{}  Median :0.0000   Median :0.0000   Median :0.0000   Median :0.00000  }
\CommentTok{\#\textgreater{}  Mean   :0.1608   Mean   :0.1391   Mean   :0.1375   Mean   :0.04385  }
\CommentTok{\#\textgreater{}  3rd Qu.:0.0000   3rd Qu.:0.0000   3rd Qu.:0.0000   3rd Qu.:0.00000  }
\CommentTok{\#\textgreater{}  Max.   :1.0000   Max.   :1.0000   Max.   :1.0000   Max.   :1.00000  }
\CommentTok{\#\textgreater{}                                                                      }
\CommentTok{\#\textgreater{}  temporada\_primavera temporada\_verano temporada\_otoño }
\CommentTok{\#\textgreater{}  Min.   :0.0000      Min.   :0.0000   Min.   :0.0000  }
\CommentTok{\#\textgreater{}  1st Qu.:0.0000      1st Qu.:0.0000   1st Qu.:0.0000  }
\CommentTok{\#\textgreater{}  Median :0.0000      Median :0.0000   Median :0.0000  }
\CommentTok{\#\textgreater{}  Mean   :0.2771      Mean   :0.2412   Mean   :0.2351  }
\CommentTok{\#\textgreater{}  3rd Qu.:1.0000      3rd Qu.:0.0000   3rd Qu.:0.0000  }
\CommentTok{\#\textgreater{}  Max.   :1.0000      Max.   :1.0000   Max.   :1.0000  }
\CommentTok{\#\textgreater{} }
\end{Highlighting}
\end{Shaded}

\subsection{Partición}\label{particiuxf3n}

Entrenamiento 2021--2024, validación 2025, como decidió el equipo (\#54,
\#55). \texttt{sima\_scores} ya trae \texttt{particion} calculada así
desde el PCA --- se reutiliza en vez de recalcularla a partir de
\texttt{fecha}, para garantizar que el PCA y los modelos de
clasificación usan exactamente la misma frontera y que ninguno vio 2025
durante el ajuste.

\begin{Shaded}
\begin{Highlighting}[]
\FunctionTok{table}\NormalTok{(sima\_disc}\SpecialCharTok{$}\NormalTok{particion)}
\CommentTok{\#\textgreater{} }
\CommentTok{\#\textgreater{} entrenamiento    validacion }
\CommentTok{\#\textgreater{}         16058          1341}

\NormalTok{sima\_train }\OtherTok{\textless{}{-}}\NormalTok{ sima\_disc }\SpecialCharTok{|\textgreater{}} \FunctionTok{filter}\NormalTok{(particion }\SpecialCharTok{==} \StringTok{"entrenamiento"}\NormalTok{)}
\NormalTok{sima\_valid }\OtherTok{\textless{}{-}}\NormalTok{ sima\_disc }\SpecialCharTok{|\textgreater{}} \FunctionTok{filter}\NormalTok{(particion }\SpecialCharTok{==} \StringTok{"validacion"}\NormalTok{)}

\FunctionTok{range}\NormalTok{(sima\_train}\SpecialCharTok{$}\NormalTok{fecha)}
\CommentTok{\#\textgreater{} [1] "2021{-}01{-}01 UTC" "2024{-}12{-}31 UTC"}
\FunctionTok{range}\NormalTok{(sima\_valid}\SpecialCharTok{$}\NormalTok{fecha)}
\CommentTok{\#\textgreater{} [1] "2025{-}01{-}01 UTC" "2025{-}06{-}30 UTC"}
\end{Highlighting}
\end{Shaded}

\textbf{Limitación de la partición} (cuantificada en
\texttt{docs/05\_estacionalidad\_ozono.pdf}): como 2025 termina el 30 de
junio, la validación es 52.3 \% primavera, 32.2 \% invierno, 15.5 \%
verano y 0 \% otoño, contra un entrenamiento repartido casi por igual
(23.1--26.0 \%). La tasa base de excedencia pasa de 29.31 \% a 41.50 \%,
12.2 puntos que vienen del calendario y no del modelo. El AUC es menos
sensible a la tasa base que la exactitud, lo que amortigua el problema
pero no lo elimina --- por eso la exactitud no se usa como métrica de
decisión más adelante.

\begin{Shaded}
\begin{Highlighting}[]
\FunctionTok{mean}\NormalTok{(sima\_train}\SpecialCharTok{$}\NormalTok{excede\_anio\_5, }\AttributeTok{na.rm =} \ConstantTok{TRUE}\NormalTok{)}
\CommentTok{\#\textgreater{} [1] 0.2940652}
\FunctionTok{mean}\NormalTok{(sima\_valid}\SpecialCharTok{$}\NormalTok{excede\_anio\_5, }\AttributeTok{na.rm =} \ConstantTok{TRUE}\NormalTok{)}
\CommentTok{\#\textgreater{} [1] 0.4150943}
\end{Highlighting}
\end{Shaded}

\subsection{Regresión logística}\label{regresiuxf3n-loguxedstica}

Respuesta \texttt{excede\_anio\_5}; predictores las tres componentes
retenidas en \#54 más los 12 dummies (zona, tipo\_dia, temporada) del
paso anterior. Se arma la fórmula a partir de los nombres de columna en
vez de escribirla a mano, para que quede ligada a los dummies que de
verdad existen en \texttt{sima\_disc} --- si cambia el catálogo de
zonas, la fórmula lo sigue sin que haya que tocarla.

\begin{Shaded}
\begin{Highlighting}[]
\NormalTok{predictores }\OtherTok{\textless{}{-}} \FunctionTok{c}\NormalTok{(}
    \StringTok{"F1"}\NormalTok{, }\StringTok{"F2"}\NormalTok{, }\StringTok{"F3"}\NormalTok{,}
    \FunctionTok{grep}\NormalTok{(}\StringTok{"\^{}zona\_|\^{}tipo\_dia\_|\^{}temporada\_"}\NormalTok{, }\FunctionTok{names}\NormalTok{(sima\_disc), }\AttributeTok{value =} \ConstantTok{TRUE}\NormalTok{)}
\NormalTok{)}
\NormalTok{formula\_log }\OtherTok{\textless{}{-}} \FunctionTok{reformulate}\NormalTok{(predictores, }\AttributeTok{response =} \StringTok{"excede\_anio\_5"}\NormalTok{)}
\NormalTok{formula\_log}
\CommentTok{\#\textgreater{} excede\_anio\_5 \textasciitilde{} F1 + F2 + F3 + zona\_Centro + zona\_Noroeste + }
\CommentTok{\#\textgreater{}     zona\_Norte + zona\_Sur + zona\_Sureste + zona\_Suroeste + tipo\_dia\_sabado + }
\CommentTok{\#\textgreater{}     tipo\_dia\_domingo + tipo\_dia\_festivo + temporada\_primavera + }
\CommentTok{\#\textgreater{}     temporada\_verano + temporada\_otoño}
\end{Highlighting}
\end{Shaded}

Ajuste sobre \texttt{sima\_train} únicamente --- \texttt{sima\_valid} no
participa del ajuste, solo de la evaluación:

\begin{Shaded}
\begin{Highlighting}[]
\NormalTok{modelo\_log }\OtherTok{\textless{}{-}} \FunctionTok{glm}\NormalTok{(formula\_log, }\AttributeTok{data =}\NormalTok{ sima\_train, }\AttributeTok{family =}\NormalTok{ binomial)}
\FunctionTok{summary}\NormalTok{(modelo\_log)}
\CommentTok{\#\textgreater{} }
\CommentTok{\#\textgreater{} Call:}
\CommentTok{\#\textgreater{} glm(formula = formula\_log, family = binomial, data = sima\_train)}
\CommentTok{\#\textgreater{} }
\CommentTok{\#\textgreater{} Coefficients:}
\CommentTok{\#\textgreater{}                     Estimate Std. Error z value Pr(\textgreater{}|z|)    }
\CommentTok{\#\textgreater{} (Intercept)         {-}2.74134    0.08013 {-}34.212  \textless{} 2e{-}16 ***}
\CommentTok{\#\textgreater{} F1                   0.65329    0.02642  24.723  \textless{} 2e{-}16 ***}
\CommentTok{\#\textgreater{} F2                  {-}0.10279    0.02464  {-}4.171 3.03e{-}05 ***}
\CommentTok{\#\textgreater{} F3                   1.02058    0.02927  34.869  \textless{} 2e{-}16 ***}
\CommentTok{\#\textgreater{} zona\_Centro          0.99142    0.08347  11.878  \textless{} 2e{-}16 ***}
\CommentTok{\#\textgreater{} zona\_Noroeste        1.32235    0.07329  18.044  \textless{} 2e{-}16 ***}
\CommentTok{\#\textgreater{} zona\_Norte           0.55689    0.08152   6.831 8.42e{-}12 ***}
\CommentTok{\#\textgreater{} zona\_Sur             1.06285    0.08744  12.156  \textless{} 2e{-}16 ***}
\CommentTok{\#\textgreater{} zona\_Sureste         0.61570    0.06913   8.907  \textless{} 2e{-}16 ***}
\CommentTok{\#\textgreater{} zona\_Suroeste        1.24327    0.07254  17.139  \textless{} 2e{-}16 ***}
\CommentTok{\#\textgreater{} tipo\_dia\_sabado      0.12578    0.05986   2.101   0.0356 *  }
\CommentTok{\#\textgreater{} tipo\_dia\_domingo     0.36157    0.06108   5.919 3.24e{-}09 ***}
\CommentTok{\#\textgreater{} tipo\_dia\_festivo     0.16145    0.10571   1.527   0.1267    }
\CommentTok{\#\textgreater{} temporada\_primavera  1.25520    0.07229  17.364  \textless{} 2e{-}16 ***}
\CommentTok{\#\textgreater{} temporada\_verano     0.41434    0.08613   4.811 1.50e{-}06 ***}
\CommentTok{\#\textgreater{} temporada\_otoño      1.10732    0.06991  15.840  \textless{} 2e{-}16 ***}
\CommentTok{\#\textgreater{} {-}{-}{-}}
\CommentTok{\#\textgreater{} Signif. codes:  0 \textquotesingle{}***\textquotesingle{} 0.001 \textquotesingle{}**\textquotesingle{} 0.01 \textquotesingle{}*\textquotesingle{} 0.05 \textquotesingle{}.\textquotesingle{} 0.1 \textquotesingle{} \textquotesingle{} 1}
\CommentTok{\#\textgreater{} }
\CommentTok{\#\textgreater{} (Dispersion parameter for binomial family taken to be 1)}
\CommentTok{\#\textgreater{} }
\CommentTok{\#\textgreater{}     Null deviance: 19025  on 15703  degrees of freedom}
\CommentTok{\#\textgreater{} Residual deviance: 15102  on 15688  degrees of freedom}
\CommentTok{\#\textgreater{}   (354 observations deleted due to missingness)}
\CommentTok{\#\textgreater{} AIC: 15134}
\CommentTok{\#\textgreater{} }
\CommentTok{\#\textgreater{} Number of Fisher Scoring iterations: 5}
\end{Highlighting}
\end{Shaded}

\subsubsection{Razones de momios}\label{razones-de-momios}

La tabla que cuenta la historia (issue \#55): \texttt{exp(coeficiente)}
convierte la escala logit a razón de momios, e \texttt{IC\ 95\%} viene
de \texttt{broom::tidy} con \texttt{conf.int\ =\ TRUE} (Wald, no perfil
de verosimilitud --- más rápido y suficiente para reportar signo y
magnitud).

\begin{Shaded}
\begin{Highlighting}[]
\FunctionTok{library}\NormalTok{(broom)}

\NormalTok{momios\_log }\OtherTok{\textless{}{-}} \FunctionTok{tidy}\NormalTok{(modelo\_log, }\AttributeTok{exponentiate =} \ConstantTok{TRUE}\NormalTok{, }\AttributeTok{conf.int =} \ConstantTok{TRUE}\NormalTok{) }\SpecialCharTok{|\textgreater{}}
    \FunctionTok{filter}\NormalTok{(term }\SpecialCharTok{!=} \StringTok{"(Intercept)"}\NormalTok{)}
\NormalTok{momios\_log}
\CommentTok{\#\textgreater{} \# A tibble: 15 x 7}
\CommentTok{\#\textgreater{}    term                estimate std.error statistic   p.value conf.low conf.high}
\CommentTok{\#\textgreater{}    \textless{}chr\textgreater{}                  \textless{}dbl\textgreater{}     \textless{}dbl\textgreater{}     \textless{}dbl\textgreater{}     \textless{}dbl\textgreater{}    \textless{}dbl\textgreater{}     \textless{}dbl\textgreater{}}
\CommentTok{\#\textgreater{}  1 F1                     1.92     0.0264     24.7  6.07e{-}135    1.83      2.02 }
\CommentTok{\#\textgreater{}  2 F2                     0.902    0.0246     {-}4.17 3.03e{-}  5    0.860     0.947}
\CommentTok{\#\textgreater{}  3 F3                     2.77     0.0293     34.9  2.16e{-}266    2.62      2.94 }
\CommentTok{\#\textgreater{}  4 zona\_Centro            2.70     0.0835     11.9  1.55e{-} 32    2.29      3.17 }
\CommentTok{\#\textgreater{}  5 zona\_Noroeste          3.75     0.0733     18.0  8.83e{-} 73    3.25      4.33 }
\CommentTok{\#\textgreater{}  6 zona\_Norte             1.75     0.0815      6.83 8.42e{-} 12    1.49      2.05 }
\CommentTok{\#\textgreater{}  7 zona\_Sur               2.89     0.0874     12.2  5.34e{-} 34    2.44      3.44 }
\CommentTok{\#\textgreater{}  8 zona\_Sureste           1.85     0.0691      8.91 5.24e{-} 19    1.62      2.12 }
\CommentTok{\#\textgreater{}  9 zona\_Suroeste          3.47     0.0725     17.1  7.58e{-} 66    3.01      4.00 }
\CommentTok{\#\textgreater{} 10 tipo\_dia\_sabado        1.13     0.0599      2.10 3.56e{-}  2    1.01      1.27 }
\CommentTok{\#\textgreater{} 11 tipo\_dia\_domingo       1.44     0.0611      5.92 3.24e{-}  9    1.27      1.62 }
\CommentTok{\#\textgreater{} 12 tipo\_dia\_festivo       1.18     0.106       1.53 1.27e{-}  1    0.953     1.44 }
\CommentTok{\#\textgreater{} 13 temporada\_primavera    3.51     0.0723     17.4  1.56e{-} 67    3.05      4.04 }
\CommentTok{\#\textgreater{} 14 temporada\_verano       1.51     0.0861      4.81 1.50e{-}  6    1.28      1.79 }
\CommentTok{\#\textgreater{} 15 temporada\_otoño        3.03     0.0699     15.8  1.64e{-} 56    2.64      3.47}
\end{Highlighting}
\end{Shaded}

\subsubsection{Razones de momios por
factor}\label{razones-de-momios-por-factor}

Los tres factores comunes, con su nombre interpretativo (no solo
``F1/F2/F3''): sin nombre, una razón de momios sobre un eje latente es
ilegible.

Ojo con el signo del segundo. \textbf{La ventilación apunta al revés que
el eje de estancamiento} que usaba la versión con componentes: valores
altos significan viento fuerte y buena dispersión, así que la razón de
momios esperada es \emph{menor} que 1. Un momio por debajo de 1 aquí no
contradice al análisis anterior, dice lo mismo con el eje orientado en
sentido contrario.

\begin{Shaded}
\begin{Highlighting}[]
\NormalTok{nombres\_factor }\OtherTok{\textless{}{-}} \FunctionTok{c}\NormalTok{(}
    \AttributeTok{F1 =} \StringTok{"Combustión primaria"}\NormalTok{,}
    \AttributeTok{F2 =} \StringTok{"Ventilación"}\NormalTok{,}
    \AttributeTok{F3 =} \StringTok{"Forzamiento fotoquímico"}
\NormalTok{)}

\NormalTok{momios\_factor }\OtherTok{\textless{}{-}}\NormalTok{ momios\_log }\SpecialCharTok{|\textgreater{}}
    \FunctionTok{filter}\NormalTok{(term }\SpecialCharTok{\%in\%} \FunctionTok{names}\NormalTok{(nombres\_factor)) }\SpecialCharTok{|\textgreater{}}
    \FunctionTok{mutate}\NormalTok{(}\AttributeTok{eje =}\NormalTok{ nombres\_factor[term], }\AttributeTok{.before =} \DecValTok{1}\NormalTok{) }\SpecialCharTok{|\textgreater{}}
    \FunctionTok{select}\NormalTok{(eje,}
        \AttributeTok{termino =}\NormalTok{ term, }\AttributeTok{momios =}\NormalTok{ estimate,}
        \AttributeTok{ic\_inf =}\NormalTok{ conf.low, }\AttributeTok{ic\_sup =}\NormalTok{ conf.high, p.value}
\NormalTok{    )}
\NormalTok{momios\_factor}
\CommentTok{\#\textgreater{} \# A tibble: 3 x 6}
\CommentTok{\#\textgreater{}   eje                     termino momios ic\_inf ic\_sup   p.value}
\CommentTok{\#\textgreater{}   \textless{}chr\textgreater{}                   \textless{}chr\textgreater{}    \textless{}dbl\textgreater{}  \textless{}dbl\textgreater{}  \textless{}dbl\textgreater{}     \textless{}dbl\textgreater{}}
\CommentTok{\#\textgreater{} 1 Combustión primaria     F1       1.92   1.83   2.02  6.07e{-}135}
\CommentTok{\#\textgreater{} 2 Ventilación             F2       0.902  0.860  0.947 3.03e{-}  5}
\CommentTok{\#\textgreater{} 3 Forzamiento fotoquímico F3       2.77   2.62   2.94  2.16e{-}266}
\end{Highlighting}
\end{Shaded}

Cada razón de momios es \textbf{por desviación estándar} del factor, no
por unidad arbitraria, gracias al reescalado de la sección anterior ---
y es eso, y no la construcción de los puntajes, lo que hace comparables
los tres momios entre sí.

\begin{Shaded}
\begin{Highlighting}[]
\FunctionTok{library}\NormalTok{(ggplot2)}

\NormalTok{fig\_momios }\OtherTok{\textless{}{-}} \FunctionTok{ggplot}\NormalTok{(momios\_factor, }\FunctionTok{aes}\NormalTok{(}\AttributeTok{x =}\NormalTok{ momios, }\AttributeTok{y =} \FunctionTok{reorder}\NormalTok{(eje, momios))) }\SpecialCharTok{+}
    \FunctionTok{geom\_vline}\NormalTok{(}\AttributeTok{xintercept =} \DecValTok{1}\NormalTok{, }\AttributeTok{linetype =} \StringTok{"dashed"}\NormalTok{, }\AttributeTok{colour =} \StringTok{"grey60"}\NormalTok{) }\SpecialCharTok{+}
    \FunctionTok{geom\_pointrange}\NormalTok{(}\FunctionTok{aes}\NormalTok{(}\AttributeTok{xmin =}\NormalTok{ ic\_inf, }\AttributeTok{xmax =}\NormalTok{ ic\_sup), }\AttributeTok{colour =} \StringTok{"\#2c7fb8"}\NormalTok{, }\AttributeTok{linewidth =} \FloatTok{0.8}\NormalTok{) }\SpecialCharTok{+}
    \FunctionTok{geom\_text}\NormalTok{(}\FunctionTok{aes}\NormalTok{(}\AttributeTok{label =} \FunctionTok{sprintf}\NormalTok{(}\StringTok{"\%.2f"}\NormalTok{, momios)), }\AttributeTok{vjust =} \SpecialCharTok{{-}}\DecValTok{1}\NormalTok{, }\AttributeTok{size =} \FloatTok{3.5}\NormalTok{) }\SpecialCharTok{+}
    \FunctionTok{labs}\NormalTok{(}
        \AttributeTok{title =} \StringTok{"Razones de momios por componente"}\NormalTok{,}
        \AttributeTok{subtitle =} \StringTok{"Excedencia MDA8 \textgreater{} 51 ppb — IC 95 \%"}\NormalTok{,}
        \AttributeTok{x =} \StringTok{"Razón de momios"}\NormalTok{, }\AttributeTok{y =} \ConstantTok{NULL}
\NormalTok{    ) }\SpecialCharTok{+}
    \FunctionTok{theme\_minimal}\NormalTok{()}

\NormalTok{fig\_momios}
\end{Highlighting}
\end{Shaded}

\pandocbounded{\includegraphics[keepaspectratio]{analisis_discriminante_files/figure-pdf/unnamed-chunk-12-1.pdf}}

\begin{Shaded}
\begin{Highlighting}[]
\FunctionTok{ggsave}\NormalTok{(}\StringTok{"../reports/figuras/momios\_componentes.pdf"}\NormalTok{, fig\_momios,}
    \AttributeTok{width =} \DecValTok{6}\NormalTok{, }\AttributeTok{height =} \FloatTok{3.6}\NormalTok{, }\AttributeTok{units =} \StringTok{"in"}
\NormalTok{)}
\end{Highlighting}
\end{Shaded}

Los tres quedan por encima de 1 con el intervalo lejos de esa línea:
subir una desviación estándar en cualquiera de los tres ejes se asocia
con mayores momios de excedencia, controlando por zona y calendario. El
de mayor efecto es el forzamiento fotoquímico, seguido de la combustión
primaria --- consistente con que el ozono es un contaminante secundario
que necesita tanto precursores como las condiciones para formarse
fotoquímicamente, no solo lo primero.

\subsubsection{Evaluación en validación: AUC, sensibilidad,
especificidad}\label{evaluaciuxf3n-en-validaciuxf3n-auc-sensibilidad-especificidad}

\textbf{Nunca exactitud} --- con 41.50 \% de excedencia base en
validación, un clasificador constante ya acierta bien sin distinguir
nada.

\begin{Shaded}
\begin{Highlighting}[]
\FunctionTok{library}\NormalTok{(pROC)}

\NormalTok{p\_valid }\OtherTok{\textless{}{-}} \FunctionTok{predict}\NormalTok{(modelo\_log, }\AttributeTok{newdata =}\NormalTok{ sima\_valid, }\AttributeTok{type =} \StringTok{"response"}\NormalTok{)}
\NormalTok{roc\_log }\OtherTok{\textless{}{-}} \FunctionTok{roc}\NormalTok{(sima\_valid}\SpecialCharTok{$}\NormalTok{excede\_anio\_5, p\_valid, }\AttributeTok{quiet =} \ConstantTok{TRUE}\NormalTok{)}

\FunctionTok{auc}\NormalTok{(roc\_log)}
\CommentTok{\#\textgreater{} Area under the curve: 0.8995}
\end{Highlighting}
\end{Shaded}

Criterio de logro del issue: AUC ≥ 0.75 en validación. Umbral de
clasificación para sensibilidad/especificidad y matriz de confusión ---
el que maximiza sensibilidad + especificidad (Youden), declarado
explícitamente en vez de asumir 0.5:

\begin{Shaded}
\begin{Highlighting}[]
\NormalTok{umbral\_log }\OtherTok{\textless{}{-}} \FunctionTok{coords}\NormalTok{(roc\_log, }\StringTok{"best"}\NormalTok{, }\AttributeTok{best.method =} \StringTok{"youden"}\NormalTok{, }\AttributeTok{ret =} \StringTok{"threshold"}\NormalTok{)}
\NormalTok{umbral\_log }\OtherTok{\textless{}{-}} \FunctionTok{as.numeric}\NormalTok{(umbral\_log)}
\NormalTok{umbral\_log}
\CommentTok{\#\textgreater{} [1] 0.2733512}

\NormalTok{pred\_log }\OtherTok{\textless{}{-}} \FunctionTok{factor}\NormalTok{(p\_valid }\SpecialCharTok{\textgreater{}=}\NormalTok{ umbral\_log,}
    \AttributeTok{levels =} \FunctionTok{c}\NormalTok{(}\ConstantTok{FALSE}\NormalTok{, }\ConstantTok{TRUE}\NormalTok{),}
    \AttributeTok{labels =} \FunctionTok{c}\NormalTok{(}\StringTok{"no\_excede"}\NormalTok{, }\StringTok{"excede"}\NormalTok{)}
\NormalTok{)}
\NormalTok{obs\_log }\OtherTok{\textless{}{-}} \FunctionTok{factor}\NormalTok{(sima\_valid}\SpecialCharTok{$}\NormalTok{excede\_anio\_5,}
    \AttributeTok{levels =} \FunctionTok{c}\NormalTok{(}\DecValTok{0}\NormalTok{, }\DecValTok{1}\NormalTok{),}
    \AttributeTok{labels =} \FunctionTok{c}\NormalTok{(}\StringTok{"no\_excede"}\NormalTok{, }\StringTok{"excede"}\NormalTok{)}
\NormalTok{)}

\NormalTok{matriz\_log }\OtherTok{\textless{}{-}} \FunctionTok{table}\NormalTok{(}\AttributeTok{observado =}\NormalTok{ obs\_log, }\AttributeTok{predicho =}\NormalTok{ pred\_log)}
\NormalTok{matriz\_log}
\CommentTok{\#\textgreater{}            predicho}
\CommentTok{\#\textgreater{} observado   no\_excede excede}
\CommentTok{\#\textgreater{}   no\_excede       635    140}
\CommentTok{\#\textgreater{}   excede           82    468}

\FunctionTok{coords}\NormalTok{(roc\_log, umbral\_log, }\AttributeTok{ret =} \FunctionTok{c}\NormalTok{(}\StringTok{"sensitivity"}\NormalTok{, }\StringTok{"specificity"}\NormalTok{))}
\CommentTok{\#\textgreater{}   sensitivity specificity}
\CommentTok{\#\textgreater{} 1   0.8509091   0.8193548}
\end{Highlighting}
\end{Shaded}

La curva ROC y la matriz de confusión no se grafican aquí: se construyen
más abajo con los dos modelos en una sola figura
(\texttt{roc\_modelos.pdf} y \texttt{matriz\_confusion.pdf}), que es
como las pide el issue \#55. Las versiones provisionales solo del
logístico (\texttt{roc\_log.pdf}, \texttt{matriz\_confusion\_log.pdf})
quedaron retiradas al existir la comparativa.

\subsubsection{Análisis de umbral}\label{anuxe1lisis-de-umbral}

Cada punto de la curva ROC es un umbral distinto; esta tabla lo hace
explícito para un puñado de candidatos, en vez de comprometerse solo con
el de Youden. Se incluye exactitud a propósito, para mostrar por qué
\textbf{no} se usa como criterio: sube con el umbral más alto (clasifica
casi todo como ``no excede'', la clase mayoritaria) mientras la
sensibilidad se desploma --- exactamente el efecto que el checklist de
la issue \#55 advierte.

\begin{Shaded}
\begin{Highlighting}[]
\NormalTok{candidatos }\OtherTok{\textless{}{-}} \FunctionTok{c}\NormalTok{(}
    \AttributeTok{youden =}\NormalTok{ umbral\_log, }\StringTok{\textasciigrave{}}\AttributeTok{0.5}\StringTok{\textasciigrave{}} \OtherTok{=} \FloatTok{0.5}\NormalTok{,}
    \AttributeTok{tasa\_base\_train =} \FunctionTok{mean}\NormalTok{(sima\_train}\SpecialCharTok{$}\NormalTok{excede\_anio\_5, }\AttributeTok{na.rm =} \ConstantTok{TRUE}\NormalTok{),}
    \AttributeTok{tasa\_base\_valid =} \FunctionTok{mean}\NormalTok{(sima\_valid}\SpecialCharTok{$}\NormalTok{excede\_anio\_5, }\AttributeTok{na.rm =} \ConstantTok{TRUE}\NormalTok{)}
\NormalTok{)}

\NormalTok{tabla\_umbrales }\OtherTok{\textless{}{-}} \FunctionTok{coords}\NormalTok{(roc\_log, candidatos,}
    \AttributeTok{ret =} \FunctionTok{c}\NormalTok{(}\StringTok{"threshold"}\NormalTok{, }\StringTok{"sensitivity"}\NormalTok{, }\StringTok{"specificity"}\NormalTok{, }\StringTok{"accuracy"}\NormalTok{)}
\NormalTok{) }\SpecialCharTok{|\textgreater{}}
    \FunctionTok{mutate}\NormalTok{(}\AttributeTok{criterio =} \FunctionTok{names}\NormalTok{(candidatos), }\AttributeTok{.before =} \DecValTok{1}\NormalTok{)}

\NormalTok{tabla\_umbrales}
\CommentTok{\#\textgreater{}                        criterio threshold sensitivity specificity  accuracy}
\CommentTok{\#\textgreater{} youden                   youden 0.2733512   0.8509091   0.8193548 0.8324528}
\CommentTok{\#\textgreater{} 0.5                         0.5 0.5000000   0.4963636   0.9509677 0.7622642}
\CommentTok{\#\textgreater{} tasa\_base\_train tasa\_base\_train 0.2940652   0.8181818   0.8438710 0.8332075}
\CommentTok{\#\textgreater{} tasa\_base\_valid tasa\_base\_valid 0.4150943   0.6309091   0.9135484 0.7962264}
\end{Highlighting}
\end{Shaded}

Sensibilidad y especificidad se mueven en direcciones opuestas al mover
el umbral --- no hay un punto que maximice ambas a la vez, por eso
Youden (maximizar su suma) es un criterio razonable por defecto y no un
óptimo universal: si el costo de una excedencia no detectada fuera mucho
mayor que el de una falsa alarma, convendría un umbral más bajo que
sacrifique especificidad por sensibilidad.

\subsection{Análisis discriminante
lineal}\label{anuxe1lisis-discriminante-lineal}

La segunda técnica de \texttt{tab:tecnicas}
(\texttt{reports/secciones/objetivo\_tecnicas.tex}) sobre la misma
matriz de diseño y la misma respuesta que la logística. El contraste es
lo que interesa: el discriminante sí supone normalidad multivariada
dentro de cada clase y homocedasticidad entre clases, supuestos que la
logística no necesita. Si ambos coinciden, el resultado no depende de
esos supuestos; si difieren, la diferencia se lee contra ellos.

\subsubsection{Verificación de
supuestos}\label{verificaciuxf3n-de-supuestos}

Box-Cox (\#52) dejó los predictores originales con
\texttt{\textbar{}sesgo\textbar{}\ \textless{}\ 0.5}. Los puntajes
factoriales son combinaciones lineales de esos predictores ---los pesos
de Bartlett de \texttt{af\_pesos.csv}--- así que la aproximación debería
trasladarse. Se verifica sobre los puntajes, que es donde el supuesto
opera, y por clase, que es como lo exige el discriminante:

\begin{Shaded}
\begin{Highlighting}[]
\NormalTok{sesgo }\OtherTok{\textless{}{-}} \ControlFlowTok{function}\NormalTok{(x) \{}
\NormalTok{    x }\OtherTok{\textless{}{-}}\NormalTok{ x[}\SpecialCharTok{!}\FunctionTok{is.na}\NormalTok{(x)]}
    \FunctionTok{sum}\NormalTok{((x }\SpecialCharTok{{-}} \FunctionTok{mean}\NormalTok{(x))}\SpecialCharTok{\^{}}\DecValTok{3}\NormalTok{) }\SpecialCharTok{/}\NormalTok{ (}\FunctionTok{length}\NormalTok{(x) }\SpecialCharTok{*} \FunctionTok{sd}\NormalTok{(x)}\SpecialCharTok{\^{}}\DecValTok{3}\NormalTok{)}
\NormalTok{\}}

\NormalTok{sima\_train }\SpecialCharTok{|\textgreater{}}
    \FunctionTok{filter}\NormalTok{(}\SpecialCharTok{!}\FunctionTok{is.na}\NormalTok{(excede\_anio\_5)) }\SpecialCharTok{|\textgreater{}}
    \FunctionTok{group\_by}\NormalTok{(excede\_anio\_5) }\SpecialCharTok{|\textgreater{}}
    \FunctionTok{summarise}\NormalTok{(}\FunctionTok{across}\NormalTok{(}\FunctionTok{c}\NormalTok{(F1, F2, F3), }\SpecialCharTok{\textasciitilde{}} \FunctionTok{round}\NormalTok{(}\FunctionTok{sesgo}\NormalTok{(.x), }\DecValTok{3}\NormalTok{)), }\AttributeTok{n =} \FunctionTok{n}\NormalTok{())}
\CommentTok{\#\textgreater{} \# A tibble: 2 x 5}
\CommentTok{\#\textgreater{}   excede\_anio\_5    F1    F2     F3     n}
\CommentTok{\#\textgreater{}           \textless{}dbl\textgreater{} \textless{}dbl\textgreater{} \textless{}dbl\textgreater{}  \textless{}dbl\textgreater{} \textless{}int\textgreater{}}
\CommentTok{\#\textgreater{} 1             0 0.244 0.391 {-}0.06  11086}
\CommentTok{\#\textgreater{} 2             1 0.051 0.508  0.072  4618}
\end{Highlighting}
\end{Shaded}

Los seis sesgos quedan por debajo de 0.51 en valor absoluto, así que la
aproximación a normalidad se sostiene sobre los puntajes.

La homocedasticidad es otra historia. No se usa la M de Box: con más de
15,000 observaciones rechaza cualquier diferencia por pequeña que sea,
así que su \(p\)-valor no informa. Se comparan las matrices de
covarianza directamente:

\begin{Shaded}
\begin{Highlighting}[]
\NormalTok{cov\_por\_clase }\OtherTok{\textless{}{-}} \ControlFlowTok{function}\NormalTok{(y) \{}
\NormalTok{    sima\_train }\SpecialCharTok{|\textgreater{}}
        \FunctionTok{filter}\NormalTok{(excede\_anio\_5 }\SpecialCharTok{==}\NormalTok{ y) }\SpecialCharTok{|\textgreater{}}
\NormalTok{        dplyr}\SpecialCharTok{::}\FunctionTok{select}\NormalTok{(F1, F2, F3) }\SpecialCharTok{|\textgreater{}}
        \FunctionTok{cov}\NormalTok{()}
\NormalTok{\}}

\NormalTok{S0 }\OtherTok{\textless{}{-}} \FunctionTok{cov\_por\_clase}\NormalTok{(}\DecValTok{0}\NormalTok{)}
\NormalTok{S1 }\OtherTok{\textless{}{-}} \FunctionTok{cov\_por\_clase}\NormalTok{(}\DecValTok{1}\NormalTok{)}

\FunctionTok{round}\NormalTok{(}\FunctionTok{sqrt}\NormalTok{(}\FunctionTok{diag}\NormalTok{(S1)) }\SpecialCharTok{/} \FunctionTok{sqrt}\NormalTok{(}\FunctionTok{diag}\NormalTok{(S0)), }\DecValTok{3}\NormalTok{) }\CommentTok{\# razón de desviaciones estándar}
\CommentTok{\#\textgreater{}    F1    F2    F3 }
\CommentTok{\#\textgreater{} 0.895 0.888 0.746}
\FunctionTok{c}\NormalTok{(}\AttributeTok{det\_no\_excede =} \FunctionTok{det}\NormalTok{(S0), }\AttributeTok{det\_excede =} \FunctionTok{det}\NormalTok{(S1))}
\CommentTok{\#\textgreater{} det\_no\_excede    det\_excede }
\CommentTok{\#\textgreater{}     1.0771349     0.3485807}
\end{Highlighting}
\end{Shaded}

Las razones son 0.895, 0.888 y 0.746: los días de excedencia son
\textbf{menos dispersos} que los de no excedencia, sobre todo en el eje
fotoquímico (\texttt{F3}), y la varianza generalizada cae de 1.077 a
0.349. El supuesto de covarianzas iguales no se cumple estrictamente. Se
declara aquí y se retoma al comparar los dos modelos, en vez de
sustituirlo por un discriminante cuadrático: la issue pide el lineal
como contraste de la logística, y cambiar de técnica a medio camino
cambiaría lo que se está contrastando.

Un segundo desvío, más grosero: los 12 dummies son binarios, y una
indicadora no es normal bajo ninguna transformación. El discriminante
lineal es conocido por tolerarlo --- el clasificador resultante sigue
siendo una frontera lineal razonable --- pero la evidencia de que aquí
lo tolera es empírica, y está en la comparación de más abajo, no en el
supuesto.

\subsubsection{Ajuste}\label{ajuste}

Misma fórmula, misma partición de entrenamiento. Las probabilidades
previas se dejan en las proporciones observadas (el comportamiento por
defecto de \texttt{MASS::lda}) en vez de fijarlas en 0.5: la tasa base
de excedencia es información del problema, no un artefacto del muestreo.

\texttt{MASS::select} no existe, pero \texttt{MASS} sí enmascara otras
funciones de \texttt{dplyr} al cargarse después; por eso los
\texttt{select} de aquí en adelante van calificados como
\texttt{dplyr::select}.

\begin{Shaded}
\begin{Highlighting}[]
\FunctionTok{library}\NormalTok{(MASS)}

\NormalTok{modelo\_lda }\OtherTok{\textless{}{-}} \FunctionTok{lda}\NormalTok{(formula\_log, }\AttributeTok{data =}\NormalTok{ sima\_train)}
\NormalTok{modelo\_lda}\SpecialCharTok{$}\NormalTok{prior}
\CommentTok{\#\textgreater{}         0         1 }
\CommentTok{\#\textgreater{} 0.7059348 0.2940652}
\end{Highlighting}
\end{Shaded}

Coeficientes de la única función discriminante --- con dos clases hay
\(\min(k-1, p) = 1\):

\begin{Shaded}
\begin{Highlighting}[]
\NormalTok{coef\_lda }\OtherTok{\textless{}{-}} \FunctionTok{data.frame}\NormalTok{(}
    \AttributeTok{termino =} \FunctionTok{rownames}\NormalTok{(modelo\_lda}\SpecialCharTok{$}\NormalTok{scaling),}
    \AttributeTok{LD1 =} \FunctionTok{round}\NormalTok{(}\FunctionTok{as.numeric}\NormalTok{(modelo\_lda}\SpecialCharTok{$}\NormalTok{scaling), }\DecValTok{4}\NormalTok{)}
\NormalTok{)}
\NormalTok{coef\_lda}
\CommentTok{\#\textgreater{}                termino     LD1}
\CommentTok{\#\textgreater{} 1                   F1  0.5057}
\CommentTok{\#\textgreater{} 2                   F2 {-}0.1074}
\CommentTok{\#\textgreater{} 3                   F3  0.7989}
\CommentTok{\#\textgreater{} 4          zona\_Centro  0.8198}
\CommentTok{\#\textgreater{} 5        zona\_Noroeste  1.1262}
\CommentTok{\#\textgreater{} 6           zona\_Norte  0.3943}
\CommentTok{\#\textgreater{} 7             zona\_Sur  0.8247}
\CommentTok{\#\textgreater{} 8         zona\_Sureste  0.4363}
\CommentTok{\#\textgreater{} 9        zona\_Suroeste  1.0686}
\CommentTok{\#\textgreater{} 10     tipo\_dia\_sabado  0.0855}
\CommentTok{\#\textgreater{} 11    tipo\_dia\_domingo  0.2751}
\CommentTok{\#\textgreater{} 12    tipo\_dia\_festivo {-}0.0051}
\CommentTok{\#\textgreater{} 13 temporada\_primavera  1.0404}
\CommentTok{\#\textgreater{} 14    temporada\_verano  0.1545}
\CommentTok{\#\textgreater{} 15     temporada\_otoño  0.7645}
\end{Highlighting}
\end{Shaded}

El orden y el signo reproducen los de la logística: primavera y otoño
pesan más que verano, el Noroeste y el Suroeste son las zonas de mayor
riesgo, y entre los factores manda el forzamiento fotoquímico
(\texttt{F3}, 0.799) sobre la combustión primaria (\texttt{F1}, 0.506),
con la ventilación (\texttt{F2}, −0.107) en contra, que es el signo
correcto para un eje que mide dispersión. Los coeficientes no son
razones de momios y no se leen como tales; son la dirección en la que
las dos clases se separan al máximo.

\subsubsection{Evaluación en
validación}\label{evaluaciuxf3n-en-validaciuxf3n}

Se usa la probabilidad posterior de la clase ``excede'' como puntaje
continuo, para que la ROC sea comparable con la de la logística --- no
la clase que \texttt{predict} asigna por defecto, que ya trae
incorporado un umbral de 0.5.

\begin{Shaded}
\begin{Highlighting}[]
\NormalTok{p\_valid\_lda }\OtherTok{\textless{}{-}} \FunctionTok{predict}\NormalTok{(modelo\_lda, }\AttributeTok{newdata =}\NormalTok{ sima\_valid)}\SpecialCharTok{$}\NormalTok{posterior[, }\StringTok{"1"}\NormalTok{]}
\NormalTok{roc\_lda }\OtherTok{\textless{}{-}} \FunctionTok{roc}\NormalTok{(sima\_valid}\SpecialCharTok{$}\NormalTok{excede\_anio\_5, p\_valid\_lda, }\AttributeTok{quiet =} \ConstantTok{TRUE}\NormalTok{)}

\FunctionTok{auc}\NormalTok{(roc\_lda)}
\CommentTok{\#\textgreater{} Area under the curve: 0.8997}
\end{Highlighting}
\end{Shaded}

Mismo criterio de umbral que en la logística (Youden), para que la
comparación de sensibilidad y especificidad no confunda el efecto del
modelo con el de una regla de corte distinta:

\begin{Shaded}
\begin{Highlighting}[]
\NormalTok{umbral\_lda }\OtherTok{\textless{}{-}} \FunctionTok{as.numeric}\NormalTok{(}\FunctionTok{coords}\NormalTok{(roc\_lda, }\StringTok{"best"}\NormalTok{,}
    \AttributeTok{best.method =} \StringTok{"youden"}\NormalTok{,}
    \AttributeTok{ret =} \StringTok{"threshold"}
\NormalTok{))}
\NormalTok{umbral\_lda}
\CommentTok{\#\textgreater{} [1] 0.3045608}

\NormalTok{pred\_lda }\OtherTok{\textless{}{-}} \FunctionTok{factor}\NormalTok{(p\_valid\_lda }\SpecialCharTok{\textgreater{}=}\NormalTok{ umbral\_lda,}
    \AttributeTok{levels =} \FunctionTok{c}\NormalTok{(}\ConstantTok{FALSE}\NormalTok{, }\ConstantTok{TRUE}\NormalTok{),}
    \AttributeTok{labels =} \FunctionTok{c}\NormalTok{(}\StringTok{"no\_excede"}\NormalTok{, }\StringTok{"excede"}\NormalTok{)}
\NormalTok{)}
\NormalTok{matriz\_lda }\OtherTok{\textless{}{-}} \FunctionTok{table}\NormalTok{(}\AttributeTok{observado =}\NormalTok{ obs\_log, }\AttributeTok{predicho =}\NormalTok{ pred\_lda)}
\NormalTok{matriz\_lda}
\CommentTok{\#\textgreater{}            predicho}
\CommentTok{\#\textgreater{} observado   no\_excede excede}
\CommentTok{\#\textgreater{}   no\_excede       650    125}
\CommentTok{\#\textgreater{}   excede           95    455}

\FunctionTok{coords}\NormalTok{(roc\_lda, umbral\_lda, }\AttributeTok{ret =} \FunctionTok{c}\NormalTok{(}\StringTok{"sensitivity"}\NormalTok{, }\StringTok{"specificity"}\NormalTok{))}
\CommentTok{\#\textgreater{}   sensitivity specificity}
\CommentTok{\#\textgreater{} 1   0.8272727   0.8387097}
\end{Highlighting}
\end{Shaded}

\subsection{Comparación de los dos
modelos}\label{comparaciuxf3n-de-los-dos-modelos}

\begin{Shaded}
\begin{Highlighting}[]
\NormalTok{comparacion }\OtherTok{\textless{}{-}} \FunctionTok{data.frame}\NormalTok{(}
    \AttributeTok{modelo =} \FunctionTok{c}\NormalTok{(}\StringTok{"Regresión logística"}\NormalTok{, }\StringTok{"Discriminante lineal"}\NormalTok{),}
    \AttributeTok{auc =} \FunctionTok{c}\NormalTok{(}\FunctionTok{as.numeric}\NormalTok{(}\FunctionTok{auc}\NormalTok{(roc\_log)), }\FunctionTok{as.numeric}\NormalTok{(}\FunctionTok{auc}\NormalTok{(roc\_lda))),}
    \AttributeTok{umbral =} \FunctionTok{c}\NormalTok{(umbral\_log, umbral\_lda),}
    \AttributeTok{sensibilidad =} \FunctionTok{c}\NormalTok{(}
        \FunctionTok{coords}\NormalTok{(roc\_log, umbral\_log, }\AttributeTok{ret =} \StringTok{"sensitivity"}\NormalTok{)[[}\DecValTok{1}\NormalTok{]],}
        \FunctionTok{coords}\NormalTok{(roc\_lda, umbral\_lda, }\AttributeTok{ret =} \StringTok{"sensitivity"}\NormalTok{)[[}\DecValTok{1}\NormalTok{]]}
\NormalTok{    ),}
    \AttributeTok{especificidad =} \FunctionTok{c}\NormalTok{(}
        \FunctionTok{coords}\NormalTok{(roc\_log, umbral\_log, }\AttributeTok{ret =} \StringTok{"specificity"}\NormalTok{)[[}\DecValTok{1}\NormalTok{]],}
        \FunctionTok{coords}\NormalTok{(roc\_lda, umbral\_lda, }\AttributeTok{ret =} \StringTok{"specificity"}\NormalTok{)[[}\DecValTok{1}\NormalTok{]]}
\NormalTok{    )}
\NormalTok{)}
\NormalTok{comparacion}
\CommentTok{\#\textgreater{}                 modelo       auc    umbral sensibilidad especificidad}
\CommentTok{\#\textgreater{} 1  Regresión logística 0.8994604 0.2733512    0.8509091     0.8193548}
\CommentTok{\#\textgreater{} 2 Discriminante lineal 0.8996950 0.3045608    0.8272727     0.8387097}
\end{Highlighting}
\end{Shaded}

Ambos superan el criterio de logro del issue (AUC ≥ 0.75) con holgura.
La diferencia de AUC se contrasta con la prueba de DeLong, que es la
adecuada porque las dos curvas están correlacionadas --- se calculan
sobre las mismas 1,325 observaciones de validación con respuesta
observada:

\begin{Shaded}
\begin{Highlighting}[]
\FunctionTok{roc.test}\NormalTok{(roc\_log, roc\_lda, }\AttributeTok{method =} \StringTok{"delong"}\NormalTok{)}
\CommentTok{\#\textgreater{} }
\CommentTok{\#\textgreater{}  DeLong\textquotesingle{}s test for two correlated ROC curves}
\CommentTok{\#\textgreater{} }
\CommentTok{\#\textgreater{} data:  roc\_log and roc\_lda}
\CommentTok{\#\textgreater{} Z = {-}0.24807, p{-}value = 0.8041}
\CommentTok{\#\textgreater{} alternative hypothesis: true difference in AUC is not equal to 0}
\CommentTok{\#\textgreater{} 95 percent confidence interval:}
\CommentTok{\#\textgreater{}  {-}0.002088208  0.001619000}
\CommentTok{\#\textgreater{} sample estimates:}
\CommentTok{\#\textgreater{} AUC of roc1 AUC of roc2 }
\CommentTok{\#\textgreater{}   0.8994604   0.8996950}
\end{Highlighting}
\end{Shaded}

\begin{Shaded}
\begin{Highlighting}[]
\FunctionTok{cor}\NormalTok{(p\_valid, p\_valid\_lda) }\CommentTok{\# correlación de los puntajes}
\CommentTok{\#\textgreater{} [1] 0.9975212}
\FunctionTok{mean}\NormalTok{((p\_valid }\SpecialCharTok{\textgreater{}=}\NormalTok{ umbral\_log) }\SpecialCharTok{==}\NormalTok{ (p\_valid\_lda }\SpecialCharTok{\textgreater{}=}\NormalTok{ umbral\_lda)) }\CommentTok{\# concordancia}
\CommentTok{\#\textgreater{} [1] 0.9701715}
\end{Highlighting}
\end{Shaded}

Las dos técnicas son indistinguibles en estos datos: los AUC difieren en
el cuarto decimal (0.8995 contra 0.8997, DeLong \(p = 0.80\)), los
puntajes correlacionan a 0.998 y las clasificaciones coinciden en el
97.0 \% de los días de validación. Donde difieren es en el reparto de
errores, y aquí el intercambio se invierte respecto de la versión con
componentes: ahora es la \textbf{logística} la que detecta más
excedencias (sensibilidad 0.851 contra 0.827) a cambio de más falsas
alarmas (especificidad 0.819 contra 0.839). Sigue siendo un intercambio
dentro del ruido de muestreo.

\textbf{Esa coincidencia es el resultado, no un anticlímax.} La
heterocedasticidad documentada arriba es real y aun así no mueve el
desempeño: la frontera que ambos trazan es prácticamente la misma. La
lectura que sigue ---las razones de momios de la {[}sección de
momios{]}--- no depende entonces del supuesto de normalidad que la
logística no hace y el discriminante sí, que era exactamente la razón de
ajustar los dos.

\subsubsection{Curvas ROC de ambos
modelos}\label{curvas-roc-de-ambos-modelos}

Una sola figura con las dos curvas, como pide el issue \#55:

\begin{Shaded}
\begin{Highlighting}[]
\CommentTok{\# El AUC va en la leyenda y no como anotación suelta: el texto no se recorta}
\CommentTok{\# contra el panel y cada curva queda etiquetada con su propio valor.}
\NormalTok{etiqueta }\OtherTok{\textless{}{-}} \ControlFlowTok{function}\NormalTok{(nombre, r) }\FunctionTok{sprintf}\NormalTok{(}\StringTok{"\%s (AUC = \%.3f)"}\NormalTok{, nombre, }\FunctionTok{auc}\NormalTok{(r))}
\NormalTok{rocs }\OtherTok{\textless{}{-}} \FunctionTok{setNames}\NormalTok{(}
    \FunctionTok{list}\NormalTok{(roc\_log, roc\_lda),}
    \FunctionTok{c}\NormalTok{(}
        \FunctionTok{etiqueta}\NormalTok{(}\StringTok{"Regresión logística"}\NormalTok{, roc\_log),}
        \FunctionTok{etiqueta}\NormalTok{(}\StringTok{"Discriminante lineal"}\NormalTok{, roc\_lda)}
\NormalTok{    )}
\NormalTok{)}

\CommentTok{\# Las dos curvas se superponen casi por completo, así que el discriminante va}
\CommentTok{\# punteado: sin eso la de abajo queda invisible y parecería haber una sola.}
\NormalTok{fig\_roc }\OtherTok{\textless{}{-}} \FunctionTok{ggroc}\NormalTok{(rocs, }\AttributeTok{aes =} \FunctionTok{c}\NormalTok{(}\StringTok{"colour"}\NormalTok{, }\StringTok{"linetype"}\NormalTok{), }\AttributeTok{linewidth =} \FloatTok{0.8}\NormalTok{) }\SpecialCharTok{+}
    \FunctionTok{geom\_abline}\NormalTok{(}\AttributeTok{intercept =} \DecValTok{1}\NormalTok{, }\AttributeTok{slope =} \DecValTok{1}\NormalTok{, }\AttributeTok{linetype =} \StringTok{"dashed"}\NormalTok{, }\AttributeTok{colour =} \StringTok{"grey60"}\NormalTok{) }\SpecialCharTok{+}
    \FunctionTok{scale\_colour\_manual}\NormalTok{(}\AttributeTok{values =} \FunctionTok{c}\NormalTok{(}\StringTok{"\#2c7fb8"}\NormalTok{, }\StringTok{"\#d95f02"}\NormalTok{)) }\SpecialCharTok{+}
    \FunctionTok{scale\_linetype\_manual}\NormalTok{(}\AttributeTok{values =} \FunctionTok{c}\NormalTok{(}\StringTok{"solid"}\NormalTok{, }\StringTok{"22"}\NormalTok{)) }\SpecialCharTok{+}
    \FunctionTok{labs}\NormalTok{(}
        \AttributeTok{title =} \StringTok{"Curvas ROC en validación (2025)"}\NormalTok{,}
        \AttributeTok{subtitle =} \StringTok{"Excedencia MDA8 \textgreater{} 51 ppb"}\NormalTok{,}
        \AttributeTok{x =} \StringTok{"Especificidad"}\NormalTok{, }\AttributeTok{y =} \StringTok{"Sensibilidad"}\NormalTok{,}
        \AttributeTok{colour =} \ConstantTok{NULL}\NormalTok{, }\AttributeTok{linetype =} \ConstantTok{NULL}
\NormalTok{    ) }\SpecialCharTok{+}
    \FunctionTok{coord\_equal}\NormalTok{() }\SpecialCharTok{+}
    \FunctionTok{theme\_minimal}\NormalTok{() }\SpecialCharTok{+}
    \FunctionTok{theme}\NormalTok{(}
        \AttributeTok{legend.position =} \StringTok{"bottom"}\NormalTok{, }\AttributeTok{legend.direction =} \StringTok{"vertical"}\NormalTok{,}
        \AttributeTok{legend.spacing.y =} \FunctionTok{unit}\NormalTok{(}\DecValTok{0}\NormalTok{, }\StringTok{"pt"}\NormalTok{)}
\NormalTok{    )}

\NormalTok{fig\_roc}
\end{Highlighting}
\end{Shaded}

\pandocbounded{\includegraphics[keepaspectratio]{analisis_discriminante_files/figure-pdf/unnamed-chunk-25-1.pdf}}

\begin{Shaded}
\begin{Highlighting}[]
\FunctionTok{ggsave}\NormalTok{(}\StringTok{"../reports/figuras/roc\_modelos.pdf"}\NormalTok{, fig\_roc,}
    \AttributeTok{width =} \DecValTok{6}\NormalTok{, }\AttributeTok{height =} \FloatTok{4.2}\NormalTok{, }\AttributeTok{units =} \StringTok{"in"}
\NormalTok{)}
\end{Highlighting}
\end{Shaded}

Las curvas se superponen en casi todo su recorrido, que es la versión
gráfica de la prueba de DeLong.

\subsubsection{Matrices de confusión}\label{matrices-de-confusiuxf3n}

Las dos matrices al umbral de Youden de cada modelo, en la misma figura
y con la misma escala para que se comparen de un vistazo:

\begin{Shaded}
\begin{Highlighting}[]
\NormalTok{prop\_por\_fila }\OtherTok{\textless{}{-}} \ControlFlowTok{function}\NormalTok{(m, modelo) \{}
    \FunctionTok{as.data.frame}\NormalTok{(m) }\SpecialCharTok{|\textgreater{}}
        \FunctionTok{group\_by}\NormalTok{(observado) }\SpecialCharTok{|\textgreater{}}
        \FunctionTok{mutate}\NormalTok{(}\AttributeTok{prop =}\NormalTok{ Freq }\SpecialCharTok{/} \FunctionTok{sum}\NormalTok{(Freq)) }\SpecialCharTok{|\textgreater{}}
        \FunctionTok{ungroup}\NormalTok{() }\SpecialCharTok{|\textgreater{}}
        \FunctionTok{mutate}\NormalTok{(}\AttributeTok{modelo =}\NormalTok{ modelo)}
\NormalTok{\}}

\NormalTok{matrices }\OtherTok{\textless{}{-}} \FunctionTok{bind\_rows}\NormalTok{(}
    \FunctionTok{prop\_por\_fila}\NormalTok{(matriz\_log, }\FunctionTok{sprintf}\NormalTok{(}\StringTok{"Logística (umbral \%.3f)"}\NormalTok{, umbral\_log)),}
    \FunctionTok{prop\_por\_fila}\NormalTok{(matriz\_lda, }\FunctionTok{sprintf}\NormalTok{(}\StringTok{"Discriminante (umbral \%.3f)"}\NormalTok{, umbral\_lda))}
\NormalTok{) }\SpecialCharTok{|\textgreater{}}
    \CommentTok{\# bind\_rows preserva el orden, pero facet\_wrap ordena los paneles por el}
    \CommentTok{\# factor: sin fijar los niveles quedarían alfabéticos y el discriminante}
    \CommentTok{\# aparecería antes que la logística, al revés de como se leen en el texto.}
    \FunctionTok{mutate}\NormalTok{(}
        \AttributeTok{modelo =} \FunctionTok{factor}\NormalTok{(modelo, }\AttributeTok{levels =} \FunctionTok{unique}\NormalTok{(modelo)),}
        \FunctionTok{across}\NormalTok{(}
            \FunctionTok{c}\NormalTok{(observado, predicho),}
            \SpecialCharTok{\textasciitilde{}} \FunctionTok{factor}\NormalTok{(.x,}
                \AttributeTok{levels =} \FunctionTok{c}\NormalTok{(}\StringTok{"no\_excede"}\NormalTok{, }\StringTok{"excede"}\NormalTok{),}
                \AttributeTok{labels =} \FunctionTok{c}\NormalTok{(}\StringTok{"No excede"}\NormalTok{, }\StringTok{"Excede"}\NormalTok{)}
\NormalTok{            )}
\NormalTok{        )}
\NormalTok{    )}

\NormalTok{fig\_matrices }\OtherTok{\textless{}{-}} \FunctionTok{ggplot}\NormalTok{(matrices, }\FunctionTok{aes}\NormalTok{(}\AttributeTok{x =}\NormalTok{ predicho, }\AttributeTok{y =}\NormalTok{ observado, }\AttributeTok{fill =}\NormalTok{ prop)) }\SpecialCharTok{+}
    \FunctionTok{geom\_tile}\NormalTok{(}\AttributeTok{colour =} \StringTok{"white"}\NormalTok{) }\SpecialCharTok{+}
    \CommentTok{\# El texto se pinta claro sobre las celdas oscuras y oscuro sobre las claras:}
    \CommentTok{\# en blanco fijo, los conteos de las celdas de proporción baja no se leen.}
    \FunctionTok{geom\_text}\NormalTok{(}
        \FunctionTok{aes}\NormalTok{(}
            \AttributeTok{label =} \FunctionTok{sprintf}\NormalTok{(}\StringTok{"\%d}\SpecialCharTok{\textbackslash{}n}\StringTok{(\%.1f \%\%)"}\NormalTok{, Freq, prop }\SpecialCharTok{*} \DecValTok{100}\NormalTok{),}
            \AttributeTok{colour =}\NormalTok{ prop }\SpecialCharTok{\textgreater{}} \FloatTok{0.5}
\NormalTok{        ),}
        \AttributeTok{fontface =} \StringTok{"bold"}\NormalTok{, }\AttributeTok{size =} \DecValTok{3}\NormalTok{, }\AttributeTok{show.legend =} \ConstantTok{FALSE}
\NormalTok{    ) }\SpecialCharTok{+}
    \FunctionTok{facet\_wrap}\NormalTok{(}\SpecialCharTok{\textasciitilde{}}\NormalTok{modelo) }\SpecialCharTok{+}
    \FunctionTok{scale\_fill\_gradient}\NormalTok{(}
        \AttributeTok{low =} \StringTok{"\#deebf7"}\NormalTok{, }\AttributeTok{high =} \StringTok{"\#08519c"}\NormalTok{, }\AttributeTok{limits =} \FunctionTok{c}\NormalTok{(}\DecValTok{0}\NormalTok{, }\DecValTok{1}\NormalTok{),}
        \AttributeTok{guide =} \StringTok{"none"}
\NormalTok{    ) }\SpecialCharTok{+}
    \FunctionTok{scale\_colour\_manual}\NormalTok{(}\AttributeTok{values =} \FunctionTok{c}\NormalTok{(}\StringTok{\textasciigrave{}}\AttributeTok{FALSE}\StringTok{\textasciigrave{}} \OtherTok{=} \StringTok{"\#08306b"}\NormalTok{, }\StringTok{\textasciigrave{}}\AttributeTok{TRUE}\StringTok{\textasciigrave{}} \OtherTok{=} \StringTok{"white"}\NormalTok{)) }\SpecialCharTok{+}
    \FunctionTok{labs}\NormalTok{(}
        \AttributeTok{title =} \StringTok{"Matrices de confusión en validación (2025)"}\NormalTok{,}
        \AttributeTok{x =} \StringTok{"Predicho"}\NormalTok{, }\AttributeTok{y =} \StringTok{"Observado"}
\NormalTok{    ) }\SpecialCharTok{+}
    \FunctionTok{theme\_minimal}\NormalTok{()}

\NormalTok{fig\_matrices}
\end{Highlighting}
\end{Shaded}

\pandocbounded{\includegraphics[keepaspectratio]{analisis_discriminante_files/figure-pdf/unnamed-chunk-26-1.pdf}}

\begin{Shaded}
\begin{Highlighting}[]
\FunctionTok{ggsave}\NormalTok{(}\StringTok{"../reports/figuras/matriz\_confusion.pdf"}\NormalTok{, fig\_matrices,}
    \AttributeTok{width =} \DecValTok{6}\NormalTok{, }\AttributeTok{height =} \FloatTok{3.2}\NormalTok{, }\AttributeTok{units =} \StringTok{"in"}
\NormalTok{)}
\end{Highlighting}
\end{Shaded}

\subsection{Contraste con las componentes
principales}\label{contraste-con-las-componentes-principales}

El cambio de insumo hay que justificarlo con desempeño, no solo con el
argumento conceptual de que un eje con nombre es una afirmación
factorial. La pregunta concreta: ¿pierde o gana el modelo al pasar de
las tres componentes de \#54 a los tres factores comunes?

Se ajustan ambos bloques \textbf{sobre exactamente las mismas filas}
---la intersección de los dos conjuntos de puntajes, ya que la poda de
\texttt{NOX} y \texttt{RH\_min\_diurna} cambió ligeramente qué días
quedan completos--- para que la prueba de DeLong sea pareada y la
comparación no confunda desempeño con cobertura.

\begin{Shaded}
\begin{Highlighting}[]
\NormalTok{pca\_scores }\OtherTok{\textless{}{-}} \FunctionTok{read\_parquet}\NormalTok{(}\StringTok{"../data/processed/pca\_puntuaciones\_diarias.parquet"}\NormalTok{)}
\FunctionTok{attr}\NormalTok{(pca\_scores}\SpecialCharTok{$}\NormalTok{fecha, }\StringTok{"tzone"}\NormalTok{) }\OtherTok{\textless{}{-}} \StringTok{"UTC"}

\NormalTok{comp }\OtherTok{\textless{}{-}}\NormalTok{ pca\_scores }\SpecialCharTok{|\textgreater{}}
\NormalTok{    dplyr}\SpecialCharTok{::}\FunctionTok{select}\NormalTok{(estacion, fecha, PC1, PC2, PC3) }\SpecialCharTok{|\textgreater{}}
    \FunctionTok{inner\_join}\NormalTok{(sima\_disc, }\AttributeTok{by =} \FunctionTok{c}\NormalTok{(}\StringTok{"estacion"}\NormalTok{, }\StringTok{"fecha"}\NormalTok{))}

\CommentTok{\# mismo reescalado que se aplicó a los factores, con momentos de entrenamiento}
\NormalTok{tr\_c }\OtherTok{\textless{}{-}}\NormalTok{ comp}\SpecialCharTok{$}\NormalTok{particion }\SpecialCharTok{==} \StringTok{"entrenamiento"}
\ControlFlowTok{for}\NormalTok{ (v }\ControlFlowTok{in} \FunctionTok{c}\NormalTok{(}\StringTok{"PC1"}\NormalTok{, }\StringTok{"PC2"}\NormalTok{, }\StringTok{"PC3"}\NormalTok{)) \{}
\NormalTok{    comp[[v]] }\OtherTok{\textless{}{-}}\NormalTok{ (comp[[v]] }\SpecialCharTok{{-}} \FunctionTok{mean}\NormalTok{(comp[[v]][tr\_c])) }\SpecialCharTok{/} \FunctionTok{sd}\NormalTok{(comp[[v]][tr\_c])}
\NormalTok{\}}

\NormalTok{comp\_train }\OtherTok{\textless{}{-}}\NormalTok{ comp }\SpecialCharTok{|\textgreater{}} \FunctionTok{filter}\NormalTok{(particion }\SpecialCharTok{==} \StringTok{"entrenamiento"}\NormalTok{)}
\NormalTok{comp\_valid }\OtherTok{\textless{}{-}}\NormalTok{ comp }\SpecialCharTok{|\textgreater{}} \FunctionTok{filter}\NormalTok{(particion }\SpecialCharTok{==} \StringTok{"validacion"}\NormalTok{)}
\FunctionTok{c}\NormalTok{(}\AttributeTok{filas\_totales =} \FunctionTok{nrow}\NormalTok{(comp), }\AttributeTok{entrenamiento =} \FunctionTok{nrow}\NormalTok{(comp\_train),}
  \AttributeTok{validacion =} \FunctionTok{nrow}\NormalTok{(comp\_valid))}
\CommentTok{\#\textgreater{} filas\_totales entrenamiento    validacion }
\CommentTok{\#\textgreater{}         17348         16008          1340}
\end{Highlighting}
\end{Shaded}

\begin{Shaded}
\begin{Highlighting}[]
\NormalTok{dummies\_nom }\OtherTok{\textless{}{-}} \FunctionTok{grep}\NormalTok{(}\StringTok{"\^{}zona\_|\^{}tipo\_dia\_|\^{}temporada\_"}\NormalTok{, }\FunctionTok{names}\NormalTok{(comp), }\AttributeTok{value =} \ConstantTok{TRUE}\NormalTok{)}

\NormalTok{ajusta }\OtherTok{\textless{}{-}} \ControlFlowTok{function}\NormalTok{(ejes) \{}
\NormalTok{    fml }\OtherTok{\textless{}{-}} \FunctionTok{reformulate}\NormalTok{(}\FunctionTok{c}\NormalTok{(ejes, dummies\_nom), }\AttributeTok{response =} \StringTok{"excede\_anio\_5"}\NormalTok{)}
\NormalTok{    g }\OtherTok{\textless{}{-}} \FunctionTok{glm}\NormalTok{(fml, }\AttributeTok{data =}\NormalTok{ comp\_train, }\AttributeTok{family =}\NormalTok{ binomial)}
\NormalTok{    d }\OtherTok{\textless{}{-}} \FunctionTok{lda}\NormalTok{(fml, }\AttributeTok{data =}\NormalTok{ comp\_train)}
    \FunctionTok{list}\NormalTok{(}
        \AttributeTok{log =} \FunctionTok{roc}\NormalTok{(comp\_valid}\SpecialCharTok{$}\NormalTok{excede\_anio\_5,}
                  \FunctionTok{predict}\NormalTok{(g, comp\_valid, }\AttributeTok{type =} \StringTok{"response"}\NormalTok{), }\AttributeTok{quiet =} \ConstantTok{TRUE}\NormalTok{),}
        \AttributeTok{lda =} \FunctionTok{roc}\NormalTok{(comp\_valid}\SpecialCharTok{$}\NormalTok{excede\_anio\_5,}
                  \FunctionTok{predict}\NormalTok{(d, comp\_valid)}\SpecialCharTok{$}\NormalTok{posterior[, }\StringTok{"1"}\NormalTok{], }\AttributeTok{quiet =} \ConstantTok{TRUE}\NormalTok{)}
\NormalTok{    )}
\NormalTok{\}}

\NormalTok{r\_fac }\OtherTok{\textless{}{-}} \FunctionTok{ajusta}\NormalTok{(}\FunctionTok{c}\NormalTok{(}\StringTok{"F1"}\NormalTok{, }\StringTok{"F2"}\NormalTok{, }\StringTok{"F3"}\NormalTok{))}
\NormalTok{r\_pca }\OtherTok{\textless{}{-}} \FunctionTok{ajusta}\NormalTok{(}\FunctionTok{c}\NormalTok{(}\StringTok{"PC1"}\NormalTok{, }\StringTok{"PC2"}\NormalTok{, }\StringTok{"PC3"}\NormalTok{))}

\FunctionTok{data.frame}\NormalTok{(}
    \AttributeTok{insumo =} \FunctionTok{c}\NormalTok{(}\StringTok{"Factores comunes"}\NormalTok{, }\StringTok{"Componentes principales"}\NormalTok{),}
    \AttributeTok{auc\_logistica =} \FunctionTok{c}\NormalTok{(}\FunctionTok{as.numeric}\NormalTok{(}\FunctionTok{auc}\NormalTok{(r\_fac}\SpecialCharTok{$}\NormalTok{log)), }\FunctionTok{as.numeric}\NormalTok{(}\FunctionTok{auc}\NormalTok{(r\_pca}\SpecialCharTok{$}\NormalTok{log))),}
    \AttributeTok{auc\_discriminante =} \FunctionTok{c}\NormalTok{(}\FunctionTok{as.numeric}\NormalTok{(}\FunctionTok{auc}\NormalTok{(r\_fac}\SpecialCharTok{$}\NormalTok{lda)), }\FunctionTok{as.numeric}\NormalTok{(}\FunctionTok{auc}\NormalTok{(r\_pca}\SpecialCharTok{$}\NormalTok{lda)))}
\NormalTok{)}
\CommentTok{\#\textgreater{}                    insumo auc\_logistica auc\_discriminante}
\CommentTok{\#\textgreater{} 1        Factores comunes     0.9001362         0.9004252}
\CommentTok{\#\textgreater{} 2 Componentes principales     0.8927719         0.8927602}
\end{Highlighting}
\end{Shaded}

\begin{Shaded}
\begin{Highlighting}[]
\FunctionTok{roc.test}\NormalTok{(r\_fac}\SpecialCharTok{$}\NormalTok{log, r\_pca}\SpecialCharTok{$}\NormalTok{log, }\AttributeTok{method =} \StringTok{"delong"}\NormalTok{)}
\CommentTok{\#\textgreater{} }
\CommentTok{\#\textgreater{}  DeLong\textquotesingle{}s test for two correlated ROC curves}
\CommentTok{\#\textgreater{} }
\CommentTok{\#\textgreater{} data:  r\_fac$log and r\_pca$log}
\CommentTok{\#\textgreater{} Z = 2.7584, p{-}value = 0.005809}
\CommentTok{\#\textgreater{} alternative hypothesis: true difference in AUC is not equal to 0}
\CommentTok{\#\textgreater{} 95 percent confidence interval:}
\CommentTok{\#\textgreater{}  0.002131606 0.012597076}
\CommentTok{\#\textgreater{} sample estimates:}
\CommentTok{\#\textgreater{} AUC of roc1 AUC of roc2 }
\CommentTok{\#\textgreater{}   0.9001362   0.8927719}
\end{Highlighting}
\end{Shaded}

Los factores ganan, y la ganancia es real pero pequeña: \textbf{0.9001
contra 0.8928} de AUC en la logística, siete milésimas, con DeLong
\(p = 0.006\). Con 1,325 observaciones de validación la prueba detecta
una diferencia de ese tamaño, así que el resultado es que los factores
\textbf{no cuestan desempeño} y lo mejoran marginalmente. El
discriminante se comporta igual (0.9004 contra 0.8928).

Esto confirma por la vía empírica lo que el solapamiento de subespacios
ya anticipaba en \texttt{docs/09\_extraccion\_factorial.pdf}: las
correlaciones canónicas entre ambos conjuntos de puntajes valen 0.996,
0.976 y 0.804, de modo que los dos bloques cargan casi la misma
información. \textbf{El argumento para cambiar de insumo no es el
desempeño, es la interpretación} --- y el desempeño confirma que cambiar
no cuesta nada.

Donde sí hay diferencia grande es en cómo se leen los coeficientes. La
tabla siguiente pone las razones de momios por desviación estándar de
ambos bloques; las filas \textbf{no están pareadas}, son dos listas
independientes puestas lado a lado, porque los ejes de un método no
corresponden uno a uno con los del otro.

\begin{Shaded}
\begin{Highlighting}[]
\CommentTok{\# razones de momios por desviación estándar, lado a lado}
\NormalTok{momios\_de }\OtherTok{\textless{}{-}} \ControlFlowTok{function}\NormalTok{(ejes, nombres) \{}
\NormalTok{    fml }\OtherTok{\textless{}{-}} \FunctionTok{reformulate}\NormalTok{(}\FunctionTok{c}\NormalTok{(ejes, dummies\_nom), }\AttributeTok{response =} \StringTok{"excede\_anio\_5"}\NormalTok{)}
    \FunctionTok{tidy}\NormalTok{(}\FunctionTok{glm}\NormalTok{(fml, }\AttributeTok{data =}\NormalTok{ comp\_train, }\AttributeTok{family =}\NormalTok{ binomial),}
         \AttributeTok{exponentiate =} \ConstantTok{TRUE}\NormalTok{) }\SpecialCharTok{|\textgreater{}}
        \FunctionTok{filter}\NormalTok{(term }\SpecialCharTok{\%in\%}\NormalTok{ ejes) }\SpecialCharTok{|\textgreater{}}
        \FunctionTok{transmute}\NormalTok{(}\AttributeTok{eje =}\NormalTok{ nombres[term], }\AttributeTok{momios =} \FunctionTok{round}\NormalTok{(estimate, }\DecValTok{3}\NormalTok{))}
\NormalTok{\}}

\NormalTok{fac }\OtherTok{\textless{}{-}} \FunctionTok{momios\_de}\NormalTok{(}\FunctionTok{c}\NormalTok{(}\StringTok{"F1"}\NormalTok{, }\StringTok{"F2"}\NormalTok{, }\StringTok{"F3"}\NormalTok{), nombres\_factor)}
\NormalTok{pca }\OtherTok{\textless{}{-}} \FunctionTok{momios\_de}\NormalTok{(}\FunctionTok{c}\NormalTok{(}\StringTok{"PC1"}\NormalTok{, }\StringTok{"PC2"}\NormalTok{, }\StringTok{"PC3"}\NormalTok{),}
                 \FunctionTok{c}\NormalTok{(}\AttributeTok{PC1 =} \StringTok{"Carga primaria de combustión"}\NormalTok{,}
                   \AttributeTok{PC2 =} \StringTok{"Forzamiento fotoquímico"}\NormalTok{,}
                   \AttributeTok{PC3 =} \StringTok{"Estancamiento atmosférico"}\NormalTok{))}

\FunctionTok{data.frame}\NormalTok{(}
    \AttributeTok{factor =}\NormalTok{ fac}\SpecialCharTok{$}\NormalTok{eje, }\AttributeTok{momios\_af =}\NormalTok{ fac}\SpecialCharTok{$}\NormalTok{momios,}
    \AttributeTok{componente =}\NormalTok{ pca}\SpecialCharTok{$}\NormalTok{eje, }\AttributeTok{momios\_pca =}\NormalTok{ pca}\SpecialCharTok{$}\NormalTok{momios}
\NormalTok{)}
\CommentTok{\#\textgreater{}                     factor momios\_af                   componente momios\_pca}
\CommentTok{\#\textgreater{} F1     Combustión primaria     1.926 Carga primaria de combustión      2.154}
\CommentTok{\#\textgreater{} F2             Ventilación     0.902      Forzamiento fotoquímico      2.210}
\CommentTok{\#\textgreater{} F3 Forzamiento fotoquímico     2.771    Estancamiento atmosférico      1.275}
\end{Highlighting}
\end{Shaded}

Tres lecturas de esa tabla:

\textbf{Los momios del PCA que están hoy en el reporte están mal
escalados.} Por desviación estándar valen 2.154, 2.210 y 1.275, no 1.47,
1.58 y 1.22. Esas últimas son por unidad de puntuación, y como las tres
componentes tienen desviaciones estándar distintas (1.98, 1.72 y 1.23),
tampoco eran comparables entre sí como el texto afirmaba. Hay que
corregirlo en \texttt{reports/secciones/regresion\_logistica.tex} se
cambie de insumo o no.

\textbf{La ventilación sale por debajo de 1} (0.902, IC 95 \% de 0.860 a
0.947): más viento, menores momios de excedencia. Es el mismo fenómeno
que el eje de estancamiento del PCA describía con un momio de 1.275,
leído con el eje orientado al revés. Que el intervalo no cruce el 1
confirma que el efecto existe, aunque sea el más débil de los tres.

\textbf{El forzamiento fotoquímico domina con claridad} (2.771 contra
1.926 de la combustión primaria). Con las componentes esa jerarquía no
se veía: 2.210 contra 2.154 las dejaba prácticamente empatadas. La
estructura simple del análisis factorial separa mejor los dos
mecanismos, y la consecuencia sustantiva es que la condición
meteorológica pesa más que la carga de precursores --- lo cual, sin
mediciones de COV, se describe pero no se explica.

\subsection{Limitaciones}\label{limitaciones}

\textbf{Autocorrelación.} Las 15 estaciones de un mismo día comparten el
clima sinóptico y los días consecutivos correlacionan entre sí, de modo
que las 17,029 filas día-estación que entran al ajuste y la validación
no son observaciones independientes: la n efectiva está mucho más cerca
de los 1,641 días distintos que cubre la serie. Medida sobre los
puntajes factoriales, la correlación entre estaciones dentro de un mismo
día promedia 0.53 en la respuesta y llega a 0.86 en el factor
fotoquímico, mientras que la autocorrelación de la respuesta a un día de
distancia es de 0.44: la dependencia transversal es la mayor de las dos.
Los errores estándar de \texttt{momios\_log} y los \(p\)-valores de
\texttt{summary(modelo\_log)} son, por lo tanto, optimistas --- los
intervalos de confianza reales son más anchos que los reportados. No se
corrige: agregar a día-AMM devolvería la independencia temporal pero
eliminaría la dimensión espacial, que es justamente donde están los
contrastes de zona que el modelo encuentra significativos. La decisión
es declararlo. Vale la pena notar que el AUC de validación no se ve
afectado por esto: se calcula sobre 2025, que el modelo no vio, y no
depende de ningún error estándar.

\textbf{Sin COV, ninguna afirmación de mecanismo.} El SIMA no mide
compuestos orgánicos volátiles, así que los ejes se describen como
\emph{consistentes con} combustión primaria, ventilación o forzamiento
fotoquímico, nunca como una atribución de régimen NOx--COV. Una razón de
momios de 2.77 sobre \texttt{F3} dice que los días con más forzamiento
fotoquímico tienen mayores momios de excedencia; no dice cuál de los dos
precursores limita la formación de ozono, y esa pregunta no es
contestable con estos datos.

\textbf{Modelo explicativo, no de pronóstico.} Los predictores son del
mismo día que la respuesta. El modelo describe qué condiciones acompañan
a una excedencia; no anticipa el ozono de mañana con los datos de hoy, y
por eso no se validó como serie de tiempo.

\textbf{Partición estacionalmente desbalanceada.} Ya cuantificada
arriba: la validación de 2025 no tiene otoño y le sobra primavera. El
AUC amortigua el sesgo de tasa base, pero las cifras de sensibilidad y
especificidad se leen sobre una validación más cargada de días propensos
a excedencia que un año completo.

\textbf{Verificación espacial pendiente.} El reajuste por zona para ver
si las razones de momios de las componentes se mueven entre zonas está
marcado como opcional en el issue \#55 y no se realizó. Queda como
extensión natural.




\end{document}
